% \documentclass[11pt,a4paper,twoside]{article}

\usepackage[utf8]{inputenc}

\usepackage[T1,T2A]{fontenc}

\usepackage[english.ancient,russian.ancient]{babel}

\usepackage{graphicx}

\usepackage{array}

\usepackage[doi]{samgtu-bib}

\usepackage{samgtu}

\usepackage{samgtu-name}

\usepackage{mathtools}

\mathtoolsset{showonlyrefs}

\usepackage{desclist}

\usepackage{euscript}

\usepackage{mathrsfs} 

\usepackage{multicol}

\usepackage{mathrsfs}

\usepackage{cite}

\usepackage{cmap}

\usepackage{qrcode}

\allowdisplaybreaks[4]

\usepackage[nodayofweek]{datetime}

\usepackage[thinlines]{easybmat}



\renewcommand{\JournalYear}{2018}

\renewcommand{\JournalVolume}{22}

\renewcommand{\JournalNumber}{4}



\newdate{JournalDateSign}{29}{12}{2018}% Дата подписания Вестника в печать

\renewcommand{\JournalDateSign}{\displaydate{JournalDateSign}}

\newdate{JournalDatePrint}{16}{01}{2019}% Дата выхода Вестника из печати

\renewcommand{\JournalDatePrint}{\displaydate{JournalDatePrint}}

%\renewcommand{\JournalUPL}{16.56} % Усл. печ. л.

%\renewcommand{\JournalUIL}{16.53} % Уч.-изд. л.

\renewcommand{\JournalUPL}{15.24} % Усл. печ. л.

\renewcommand{\JournalUIL}{15.21} % Уч.-изд. л.

\renewcommand{\JournalRegNumber}{220/18} % Регистрационный номер

\renewcommand{\JournalOrderNumber}{2} % Номер заказа



%\renewcommand{\JournalRMonth}{Март} % Месяц выхода Вестника

%\renewcommand{\JournalEMonth}{March} % Месяц выхода Вестника

%\renewcommand{\JournalRMonth}{Июнь} % Месяц выхода Вестника

%\renewcommand{\JournalEMonth}{June} % Месяц выхода Вестника

%\renewcommand{\JournalRMonth}{Сентябрь} % Месяц выхода Вестника

%\renewcommand{\JournalEMonth}{September} % Месяц выхода Вестника

\renewcommand{\JournalRMonth}{Декабрь} % Месяц выхода Вестника

\renewcommand{\JournalEMonth}{December} % Месяц выхода Вестника



\newcommand{\totop}[1]{\raisebox{6pt}[1pt][2pt]{#1}}

\newcommand{\tottop}[1]{\raisebox{12pt}[1pt][2pt]{#1}} 

\newcommand{\totttop}[1]{\raisebox{18pt}[1pt][2pt]{#1}} 

\newcommand{\tottttop}[1]{\raisebox{24pt}[1pt][2pt]{#1}} 

\newcommand{\totttttop}[1]{\raisebox{30pt}[1pt][2pt]{#1}} 

\newcommand{\tobot}[1]{\raisebox{-6pt}[1pt][2pt]{#1}} 

\newcommand{\tobbot}[1]{\raisebox{-12pt}[1pt][2pt]{#1}} 

\newcommand{\tobbbot}[1]{\raisebox{-18pt}[1pt][2pt]{#1}}



\graphicspath{{./00PIC/}}

%

%

%

%

\newcommand{\mmMGUI}{\ruaffid{595}{Московский государственный университет им.~М.~В. Ломоносова, \\ 

Россия, 119899, Москва, Воробьёвы горы}}

\newcommand{\engmmMGUI}{\enaffid{595}{M.~V.~Lomonosov Moscow State University, \\

 Vorob'evy gory, Moscow,   119899, Russian Federation}}

%

%

%%\graphicspath{{./00PIC/}}

%

%%\samgtupubl % для генерации pdf, отдаваемого в типографию СамГТУ

%\pdfcrop % обрезка pdf для размещения на math-net.ru

%\draft % На корректуре

%%\filllastpage % Последняя страница не заполнена не полностью

%%\briefcomm % Статья является кратким сообщением

%

\vsgtu{1504}

%

\FirstPage{249}

\LastPage{270}

%%

%%

%\begin{document}





\renewcommand{\baselinestretch}{0.9}



\hfill \qrcode[hyperlink,height=.8in]{http://dx.doi.org/10.14498/vsgtu1504}

\vspace{-.8in}







\udc{517.927}

\msc{34B10, 47E05}



\rutitle{Об эффекте <<расщепления>> для многоточечных дифференциальных операторов с суммируемым потенциалом}

\rutitleshort{Об эффекте <<расщепления>> для многоточечных дифференциальных операторов \dots}





\ruauthor{С.~И.~Митрохин}{М\,и\,т\,р\,о\,х\,и\,н~С.\,И.}



\ruaffil{\mmMGUI.}

\enaffil{\engmmMGUI.}



\ruauthorsdetails{1}{

\fullauthor{\rupid{46310}{Сергей Иванович Митрохин}} 

\CAuthor % Автор-корреспондент

%\orcid{0000-000x-xxxx-xxxx} % ORCID ID автора 

\regalia{~кандидат физико-математических наук, доцент} 

\position{старший научный сотрудник} 

\dept{научно"=исследовательский вычислительный центр} 

\email[b]{mitrokhin-sergey@yandex.ru}

}



\enauthorsdetails{1}{

\fullauthor{\enpid{46310}{Sergey I. Mitrokhin}} 

\CAuthor % Автор-корреспондент

%\orcid{0000-000x-xxxx-xxxx} % ORCID ID автора 

\regalia{~Cand. Phys. \& Math. Sci., Associate Professor}

\position{Senior Researcher} 

\dept{Research Computing Center} 

\email{mitrokhin-sergey@yandex.ru}

}





\rukeywords{дифференциальный оператор, спектральный параметр, суммируемый потенциал, уравнение на собственные значения, индикаторная диаграмма, асимптотика собственных значений}



\ruabstract{Изучается дифференциальный оператор четвёртого порядка с многоточечными граничными условиями. Потенциал дифференциального оператора является суммируемой функцией на конечном отрезке. При больших значениях спектрального параметра найдена асимптотика решений дифференциального уравнения, задающего дифференциальный оператор. 

На основе изучения граничных условий  выведено уравнение на собственные значения изучаемого оператора. Параметры граничных условий подобраны таким образом, что в главном приближении уравнение на собственные значения имеет кратные корни. Автором показано, что у исследуемого оператора наблюдается эффект <<расщепления>> кратных в~главном приближении собственных значений. Выведены все серии однократных собственных значений изучаемого оператора. Изучена индикаторная диаграмма рассматриваемого оператора. Найдена асимптотика собственных значений во всех секторах индикаторной диаграммы. Полученной точности асимптотических формул достаточно для нахождения асимптотики собственных функций исследуемого дифференциального оператора.}



\entitle{On the ``splitting'' effect for multipoint differential operators with summable potential}



\enauthor{S.~I.~Mitrokhin}{M\,i\,t\,r\,o\,k\,h\,i\,n~S.~I.}





\enabstract{We study the differential operator of the fourth order with multipoint boundary conditions. The potential of the differential operator is summable function on a finite segment. For large values of spectral parameter the asymptotic behavior of solutions of differential equation which define the differential operator is found. The equation for eigenvalues of the studied operators is derived by studying the boundary conditions. The parameters of boundary conditions are selected in such a way that the main approach of the equation for eigenvalues has multiple roots. The author shows that for the studied operator the effect of ``splitting'' of multiple eigenvalues in the main approximation is observed. We derive all series of single eigenvalues of the investigated operator. The indicator diagram of the considered operator is studied. The asymptotic behavior of eigenvalues in all sectors of the indicator diagram is found. The obtained precision of the asymptotic formulas is enough for finding an asymptotics of eigenfunctions of the studied differential operator.}



\enkeywords{differential operator, spectral parameter, summable potential, the equation for the eigenvalues, the indicator diagram, the asymptotic behavior of the eigenvalues}



\newdate{mitra}{24}{08}{2016}

\date{\displaydate{mitra}}

\newdate{mitrb}{13}{05}{2017}

\revisiondate{\displaydate{mitrb}}

\newdate{mitrc}{12}{06}{2017}

\accepteddate{\displaydate{mitrc}}



\newdate{mitre}{05}{07}{2017}

\renewcommand{\JournalDataOnline}{\displaydate{mitre}}



\makerutitle



\newpage





\Section{Постановка задачи} Исследуем дифференциальный оператор четвертого порядка, задаваемый дифференциальным уравнением вида

\begin{equation}

\label{mitr:eq1}

y^{\eqref{mitr:eq4}}(x)+q(x)y(x)=\lambda  y(x), \quad 

0\leqslant x\leqslant \pi,

\end{equation}

с многоточечными граничными условиями

\begin{equation}

\label{mitr:eq2}

\begin{gathered}

y'(0)=0; \quad  

y^{\eqref{mitr:eq3}}(0)=0; \quad 

y(\pi)=\alpha_1 y\Bigl(\frac\pi3\Bigr)+\alpha_2 y\Bigl(\frac{2\pi}3\Bigr); \\

y''(\pi)=\beta_1 y''\Bigl(\frac\pi3\Bigr)+\beta_2 y'' \Bigl(\frac{2\pi}3\Bigr);  \\ 

\beta_1=-\frac{10}3- \alpha_1, \quad  

\alpha_1\in\mathbb{C}; \quad  

\beta_2=\frac{16}3-\alpha_2, \quad  \alpha_2\in\mathbb{C}.

\end{gathered}

\end{equation}





В уравнении \eqref{mitr:eq1} функция $q(x)$  "---  потенциал,  число $\lambda\in\mathbb{C}$ "--- спектральный параметр. 

Мы изучаем дифференциальный оператор \eqref{mitr:eq1}, \eqref{mitr:eq2}  в предположении, что потенциал $q(x)$ является суммируемой функцией на отрезке $[0; \pi]$:

\begin{equation}

\label{mitr:eq3}

q(x)\in L_1[0;\pi]\Leftrightarrow \left(\int^x_0 q(t) dt\right)'_x=q(x)

%\eqno\eqref{mitr:eq3}

\end{equation}

почти для всех $x$ из отрезка $[0; \pi]$.



Многоточечные дифференциальные операторы пока что полностью не изучены даже для случая гладких коэффициентов дифференциальных уравнений, задающих эти операторы. 

В~работе~\cite{mitr:bib1} автором для дифференциального оператора второго порядка с кусочно-гладкими коэффициентами был изучен эффект <<расщепления>> кратных в~главном собственных значений. 

Цель статьи "--- найти примеры и изучить дифференциальные операторы четвертого порядка с суммируемым потенциалом, у~которых наблюдается такой же эффект <<расщепления>>.



Дифференциальные операторы с суммируемыми коэффициентами начали бурно исследоваться в последнее время. 

В~работах~\cite{mitr:bib2,mitr:bib3} были исследованы спектральные свойства оператора Штурма"--~Лиувилля (второго порядка) с~суммируемым потенциалом.



В работах \cite{mitr:bib4,mitr:bib5} автор исследовал операторы четвертого и шестого порядков с~суммируемыми коэффициентами с~помощью методики, отличной от методики работ~\cite{mitr:bib2,mitr:bib3}. 

Необходимо отметить, что с~возрастанием порядка дифференциального оператора трудности его исследования увеличиваются многократно. 

И~этот факт верен и~для операторов с~негладкими коэффициентами: фактически во всех работах этой тематики (см.~[\citen{mitr:bib6}--\citen{mitr:bib10}]) рассматриваются операторы второго порядка.



Нахождение асимптотики собственных значений необходимо при исследовании свойств собственных функций, а также для вычисления регуляризованных следов дифференциальных операторов (см.~[\citen{mitr:bib11}--\citen{mitr:bib13}]). Заметим, что общая теория нахождения регуляризованных следов операторов с суммируемыми потенциалами пока не разработана, хотя  появились работы, в которых вычислены следы операторов (опять"=таки второго порядка) c потенциалом  в виде $\delta$-функции (см.~\cite{mitr:bib14,mitr:bib15}). Для операторов четвертого и выше порядков случай, когда потенциал "--- $\delta$"=функция, пока не рассматривался.



\smallskip



\Section{Асимптотика решений дифференциального уравнения (\ref{mitr:eq1}) при больших значениях $\boldsymbol\lambda$}

Введем замену $\lambda=s^4$, $s=\sqrt[4]{\lambda}$, причем для корректности дальнейших выкладок зафиксируем ту ветвь арифметического корня, для которой $\sqrt[4]{1}=+1$. Обозначим через $w_k$ $(k=1, 2, 3, 4)$ различные корни четвертой степени из единицы:

\begin{equation}

\label{mitr:eq4}

\begin{gathered}

w^4_k=1; \quad  w_k=\exp\Bigl(\frac{2\pi i}4(k-1)\Bigr), \\

k=1, 2, 3, 4; \quad w_1=1, \quad  w_2=i, \quad w_3=-1, \quad w_4=-i.

%\eqno\eqref{mitr:eq4}

\end{gathered}

\end{equation}



Для чисел $w_k$ $(k=1, 2, 3, 4)$ из \eqref{mitr:eq4} справедливы соотношения

\begin{equation}

\label{mitr:eq5}

\sum^4_{k=1}w^p_k=0, \; p=1,2,3; \quad \sum^4_{k=1}w^p_k=4, \; p=0, \; p=4.

\end{equation}



Числа $w_k$ $(k=1, 2, 3, 4)$ из \eqref{mitr:eq4}, \eqref{mitr:eq5} делят единичную окружность на четыре равные части.



Методом вариации произвольных постоянных устанавливается следующее утверждение (см.~\cite{mitr:bib4,mitr:bib5}).



\begin{lemma}[1]Решение $y(x,s)$ дифференциального уравнения \eqref{mitr:eq1} является решением интегрального уравнения Вольтерры{\rm :}

\begin{equation}

\label{mitr:eq6}

y(x,s)=\sum^4_{k=1}C_k e^{w_ksx}-\frac1{4s^3}

\sum^4_{k=1}w_k e^{w_ksx}\int^x_0 q(t) e^{-w_kst} y(t,s)dt,

\end{equation}

где $C_k$ "--- произвольные постоянные.

\end{lemma}





Проверить справедливость формулы \eqref{mitr:eq6} можно непосредственным образом с учетом свойства \eqref{mitr:eq5} а также принимая во внимание, что из свойства \eqref{mitr:eq3}  суммируемости потенциала следует формула

$$\biggl(\int^x_0q(t) e^{-w_kst}y(t,s) dt\biggr)'_x=q(x)e^{-w_ksx} y(x,s)$$

почти для всех $x$ из отрезка $[0;\pi]$.



Интегральное уравнение \eqref{mitr:eq6} решим методом последовательных приближений Пикара (как это сделано в работах   \cite{mitr:bib4} и  \cite{mitr:bib5}): 

найдем $y(t,s)$ из \eqref{mitr:eq6}, подставим это выражение в $\displaystyle \int^x_0 q(t) e^{-w_kst} y(t,s)dt$, произведем необходимые вычисления и оценки, аналогичные оценкам работ  \cite{mitr:bib4, mitr:bib5}, \cite[гл.~2]{mitr:bib16}, \cite[гл.~1]{mitr:bib17}. В результате получим следующее утверждение.



\smallskip



\begin{lemma}[2]Общее решение дифференциального уравнения \eqref{mitr:eq1} представляется в виде

\begin{equation}

\label{mitr:eq7}

y(x,s)=\sum^4_{k=1} C_k y_k(x,s); \quad

y^{(p)}(x,s)=\sum^4_{k=1} C_k y^{(p)}_k(x,s), \quad  p=1, 2, 3,

\end{equation}

где $C_k$ $(k=1, 2, 3, 4)$ "--- произвольные постоянные$,$

\begin{equation}

\label{mitr:eq8}

y_k(x,s)=e^{w_ksx}-\frac{A_{3k}(x,s)}{4s^3}+ \frac{A_{6k}(x,s)}{16s^6}+

\underline O\Bigl(\frac{e^{|\mathrm{Im} s|x}}{s^9}\Bigr), \quad k=1, 2, 3, 4;\\

\end{equation}

\begin{equation}

\label{mitr:eq10}

\begin{gathered}

\frac{y^{(p)}_k(x,s)}{s^p}=w^p_k e^{w_ksx}-\frac{A^p_{3k}(x,s)}{4s^3}+

\frac{A^p_{6k}(x,s)}{16s^6}+

\underline O\Bigl(\frac{e^{|\mathrm{Im} s|x}}{s^9}\Bigr), \\ k=1, 2, 3, 4, \quad p=1, 2, 3;

\end{gathered}

\end{equation}

\begin{equation}

\label{mitr:eq10}

\begin{gathered}

A_{3k}(x,s)=\sum^4_{n=1} w_n e^{w_nsx}\varphi_{kn}(x,s), \quad \varphi_{kn} (x,s)=

\int ^x_0 q(x) e^{(w_k-w_n)st} dt, \\ k=1,2,3,4;

\end{gathered}

\end{equation}

\begin{equation}

\label{mitr:eq11}

A^p_{3k}(x,s)=\sum^4_{n=1} w_n w^p_n e^{w_nsx}\varphi_{kn} (x,s), \quad k=1, 2, 3, 4, \quad p=1, 2, 3;

\end{equation}

\begin{equation}

\label{mitr:eq12}

\begin{gathered}

A_{6k}(x,s)=\sum^4_{n=1} w_n \biggl(\sum^4_{m=1} w_m

 e^{w_msx}\int ^x_0 q(t) \psi_{nmkn}(x,s)\biggr),  \\

\psi_{nmkn}(x,s)= \int ^x_0 q(t)

e^{(w_n-w_m)st} d\xi

\biggl(\int ^t_0q(\xi) e^{(w_k-w_n)s\xi}d\xi\biggr) dt, \\

k=1, 2, 3, 4;

\end{gathered}

\end{equation}

\begin{equation}

\label{mitr:eq13}

\begin{gathered}

A^p_{6k}(x,s)=\sum^4_{n=1} w_n \biggl(\sum^4_{m=1}w_m w^p_m e^{w_msx}

\psi_{nmkn}(x,s)\biggr), \\ k=1,2,3,4, \quad p=1, 2, 3.

\end{gathered}

\end{equation}



\end{lemma}



\smallskip



При выводе асимптотических формул \eqref{mitr:eq8}--\eqref{mitr:eq13} в момент промежуточного интегрирования мы требовали выполнения следующих начальных условий:

\begin{equation}

\label{mitr:eq14}

\begin{gathered}

A_{3k}(0,s)=A_{6k}(0,s)=0; \quad

A^p_{3k}(0,s)=A^p_{6k}(0,s)=0; \quad  y_k(0,s)=1; \\

y^{(p)}_k(0,s)=w^p_k s^p, \quad k=1,2,3,4; \quad p=1,2,3.

%\eqno\eqref{mitr:eq16}

\end{gathered}

\end{equation}





\smallskip



\Section{Изучение граничных условий (\ref{mitr:eq2})}

В силу формул \eqref{mitr:eq7}--\eqref{mitr:eq14} из первых двух граничных условий \eqref{mitr:eq2} находим:

\begin{equation*}

\begin{cases}

\displaystyle

\frac{y'(0,s)}{s}\Big|_{s\ne0} \stackrel{(\ref{mitr:eq2})}{=} 0 \stackrel{(\ref{mitr:eq7})}{\Leftrightarrow} \sum\limits^4_{k=1}C_k \frac{y'_k(0,s)}{s}=0\stackrel{(\ref{mitr:eq14})}{\Leftrightarrow} \sum\limits^4_{k=1} C_k w_k=0; \\ 

\displaystyle

\frac{y^{(3)}(0,s)}{s^3}\Big|_{s\ne0} \stackrel{(\ref{mitr:eq2})}{=} 0 \stackrel{(\ref{mitr:eq7})}{\Leftrightarrow} \sum\limits^4_{k=1} 

C_k \frac{y^{(3)}_k(0,s)}{s^3}=0\stackrel{\eqref{mitr:eq14}}{\Leftrightarrow} \sum\limits^4_{k=1} 

C_k w^3_k=0,

\end{cases}

\end{equation*}

откуда в силу соотношений \eqref{mitr:eq5} получаем: $C_3=C_1$, $C_4=C_2$. 

Поэтому при выполнении условий $y'(0)=y^{(3)}(0)=0$ общее решение дифференциального уравнения \eqref{mitr:eq1} имеет следующий вид:

\begin{equation}

\label{mitr:eq15}

\begin{gathered}

y(x,s)=C_1[y_1(x,s)+y_3(x,s)]+C_2[y_2(x,s)+y_4(x,s)],\\

y^{(p)}(x,s)=C_1[y^{(p)}_1(x,s)+y^{(p)}_3(x,s)]+C_2[y^{(p)}_2(x,s)+y^{(p)}_4(x,s)],\\

p=1,2,3,

%\eqno\eqref{mitr:eq17}

\end{gathered}

\end{equation}

где $C_1,C_2$ "--- произвольные постоянные.



Подставляя формулы \eqref{mitr:eq15} в последние два из граничных условий \eqref{mitr:eq2}, получаем:

\begin{equation}

\label{mitr:eq16}

\!\begin{cases}

\displaystyle

y(\pi,s)\stackrel{\eqref{mitr:eq2}}{=} \alpha_1y \Bigl(\frac\pi3,s \Bigr)+

\alpha_2 y \Bigl(\frac{2\pi}3,s \Bigr)

 \stackrel{\eqref{mitr:eq15}}{\Leftrightarrow}

C_1[y_1(\pi,s)+y_3(\pi,s)]+\\

\displaystyle \qquad 

 +

C_2[y_2(\pi,s)+y_4(\pi,s)]=

\alpha_1 C_1\left[y_1\Bigl(\frac\pi3,s\Bigr)+y_3\Bigl(\frac\pi3,s\Bigr)\right]+\\

\displaystyle \qquad \qquad 

+\alpha_1 C_2\left[y_2(\frac\pi3,s)+ y_4(\frac\pi3,s)\right]+ \\

\displaystyle \qquad \qquad \qquad 

\alpha_2 C_1\left[y_1 \Bigl(\frac{2\pi}3,s\Bigr)+y_3\Bigl(\frac{2\pi}3,s\Bigr)\right]+\\

\displaystyle \qquad \qquad \qquad \qquad 

+\alpha_2 C_2\left[y_2 \Bigl(\frac{2\pi}3,s\Bigr)+y_4\Bigl(\frac{2\pi}3,s\Bigr)\right],\\[6mm]

\displaystyle

\frac{y''(\pi,s)}{s^2} \stackrel{\eqref{mitr:eq2}}{=}

\beta_1\frac{y''(\pi/3,s)}{s^2}+ \beta_2\frac{y''({2\pi}/3,s)}{s^2}

\stackrel{\eqref{mitr:eq15}}{\Leftrightarrow}

C_1\left[\frac{y''_1(\pi,s)}{s^2} +\frac{y''_3(\pi,s)}{s^2}\right]+\\ 

\displaystyle \qquad 

+

C_2\left[\frac{y''_2(\pi,s)}{s^2} +\frac{y''_4(\pi,s)}{s^2}\right]=

\beta_1 C_1\left[\frac{y''_1(\pi/3,s)}{s^2} +\frac{y''_3(\pi/3,s)}{s^2}\right]+\\

\displaystyle \qquad \qquad 

+\beta_1 C_2\left[\frac{y''_2(\pi/3,s)}{s^2} +\frac{y''_4(\pi/3,s)}{s^2}\right]+ \\

\displaystyle \qquad \qquad \qquad 

+

\beta_2 C_1\left[\frac{y''_1({2\pi}/3,s)}{s^2} +\frac{y''_3({2\pi}/3,s)}{s^2}\right]+\\

\displaystyle \qquad \qquad \qquad \qquad 

+\beta_2 C_2\left[\frac{y''_2({2\pi}/3,s)}{s^2} +\frac{y''_4({2\pi}/3,s)}{s^2}\right].

\end{cases} \hspace{-15mm}

\end{equation}





Однородная система \eqref{mitr:eq16} (из двух уравнений с двумя неизвестными~$C_1$,~$C_2$) имеет ненулевые решения ($C^2_1+C^2_2 \ne0$) только в том случае, когда ее определитель равен нулю. Поэтому справедлива следующая лемма.



\smallskip



\begin{lemma}[3] Уравнение на собственные значения дифференциального оператора \eqref{mitr:eq1}, \eqref{mitr:eq2} с условием \eqref{mitr:eq3} суммируемости потенциала $q(x)$ имеет следующий вид\/{\rm:}

\begin{equation}

\label{mitr:eq17}

f(s)= \begin{vmatrix}

h_{11} (s) &h_{12} (s)\\

h_{21} (s) &h_{22}(s)

\end{vmatrix}

=0,

\end{equation}

%\allowdisplaybreaks

\begin{multline}

h_{11}(s)=y_1(\pi,s)+y_3(\pi,s)-\alpha_1\bigg[y_1\bigg(\frac\pi3,s\bigg)+

 y_3\bigg(\frac\pi3,s\bigg)\bigg]- \\-

\alpha_2\bigg[y_1\bigg(\frac{2\pi}3,s\bigg)+  y_3\bigg(\frac{2\pi}3,s\bigg)\bigg],

\end{multline}

\begin{multline}

h_{12}(s)=y_2(\pi,s)+y_4(\pi,s)-\alpha_1\bigg[y_2\bigg(\frac\pi3,s\bigg)+

 y_4\bigg(\frac\pi3,s\bigg)\bigg]-\\ -

\alpha_2\bigg[y_2\bigg(\frac{2\pi}3,s\bigg)+  y_4\bigg(\frac{2\pi}3,s\bigg)\bigg],

\end{multline}

\begin{multline}

h_{21}(s)=\frac{y''_1(\pi,s)}{s^2}+\frac{y''_3(\pi,s)}{s^2}-\beta_1\bigg[

\frac{y''_1(\pi/3,s)}{s^2}+

\frac{y''_3(\pi/3,s)}{s^2}\bigg]-\\-

\beta_2\bigg[

\frac{y''_1({2\pi}/3,s)}{s^2}+\frac{y''_3({2\pi}/3,s)}{s^2}\bigg],

\end{multline}

\begin{multline}

h_{22}(s)=\frac{y''_2(\pi,s)}{s^2}+\frac{y''_4(\pi,s)}{s^2}-\beta_1\bigg[

\frac{y''_2(\pi/3,s)}{s^2}+

\frac{y''_4(\pi/3,s)}{s^2}\bigg]-\\-

\beta_2\bigg[

\frac{y''_2({2\pi}/3,s)}{s^2}+\frac{y''_4({2\pi}/3,s)}{s^2}\bigg].

\end{multline}

%\end{equation*}



\end{lemma}

\smallskip





Подставляя формулы \eqref{mitr:eq8}--\eqref{mitr:eq13} в уравнение \eqref{mitr:eq17}, получившийся определитель разложим на сумму определителей по столбцам, в результате получим:

\begin{equation}

\label{mitr:eq18}

f(s)=f_0(s)-\frac{f_3(s)}{4s^3}+\frac{f_6(s)}{16s^6}+\underline O\Bigl(\frac1{s^9}\Bigr)=0,\\

\end{equation}

\begin{equation}

\label{mitr:eq22}

f_0(s)= \begin{vmatrix}

f_{011} (s) &f_{012} (s)\\

f_{021} (s) &f_{022}(s)

\end{vmatrix},\\ 

\end{equation}

%\allowdisplaybreaks

$$

f_{011}(s)=

e^{w_1s\pi} +e^{w_3s\pi}-\alpha_1\Big[e^{w_1s\frac\pi3}+

e^{w_3s\frac\pi3}\Big]-\alpha_2\Big[e^{w_1s\frac{2\pi}3}+

e^{w_3s\frac{2\pi}3}\Big],

$$

$$

f_{012}(s)=

e^{w_2s\pi} +e^{w_4s\pi}-\alpha_1\Big[e^{w_2s\frac\pi3}+

e^{w_4s\frac\pi3}\Big]-\alpha_2\Big[e^{w_2s\frac{2\pi}3}+

e^{w_4s\frac{2\pi}3}\Big],

$$

\begin{multline}

f_{021}(s)=

w^2_1e^{w_1s\pi} +w^2_3e^{w_3s\pi}-\beta_1

\Big[w^2_1e^{w_1s\frac\pi3}+

w^2_3e^{w_3s\frac\pi3}\Big]-\\-\beta_2\Big[w^2_1e^{w_1s\frac{2\pi}3}+

w^2_3e^{w_3s\frac{2\pi}3}\Big],

\end{multline}

\begin{multline}

f_{022}(s)=

w^2_2e^{w_2s\pi} +w^2_4 e^{w_4s\pi}-\beta_1

\Big[w^2_2 e^{w_2s\frac\pi3}+

w^2_4 e^{w_4s\frac\pi3}\Big]-\\-\beta_2\Big[w^2_2 e^{w_2s\frac{2\pi}3}+

w^2_4 e^{w_4s\frac{2\pi}3}\Big],

\end{multline}

\begin{equation}

\label{mitr:eq22}

f_3(s)=f_{3,1}(s)+f_{3,2}(s), \quad w^2_1=w^2_3=1, \quad w^2_2=w^2_4=-1,

\end{equation}

\begin{equation}

\label{mitr:eq22}

f_{3,1}(s)= \begin{vmatrix}

f_{3,1,11} (s) &f_{3,1,12} (s)\\

f_{3,1,21} (s) &f_{3,1,22}(s)

\end{vmatrix},

\end{equation}

%\allowdisplaybreaks

\begin{multline}

f_{3,1,11}(s)=A_{31}(\pi,s)+A_{33}(\pi,s)-\alpha_1\left[A_{31}\Bigl(\frac\pi3,s\Bigr)+

A_{33}\Bigl(\frac\pi3,s\Bigr)\right] -\\-\alpha_2 \left[A_{31}\Bigl(\frac{2\pi}3,s\Bigr)+A_{33}

\Bigl(\frac{2\pi}3,s\Bigr)\right],

\end{multline}

$$

f_{3,1,12}(s)=e^{w_2s\pi}+e^{w_4s\pi}-\alpha_1\Big[e^{w_2s\frac\pi3}+

e^{w_4s\frac\pi3}\Big]-\alpha_2

\Big[e^{w_2s\frac{2\pi}3}+e^{w_4s\frac{2\pi}3}\Big],

$$

\begin{multline}

f_{3,1,21}(s)=A^2_{31}(\pi,s)+A^2_{33}(\pi,s)-\beta_1\left[A^2_{31}\Bigl(\frac\pi3,s\Bigr)+

A^2_{33}\Bigl(\frac\pi3,s\Bigr)\right] -\\-\beta_2 \left[A^2_{31}\Bigl(\frac{2\pi}3,s\Bigr)+A^2_{33}

\Bigl(\frac{2\pi}3,s\Bigr)\right],

\end{multline}

$$

f_{3,1,22}(s)=-e^{w_2s\pi}-e^{w_4s\pi}+\beta_1\Big[e^{w_2s\frac\pi3}+

e^{w_4s\frac\pi3}\Big]+\beta_2

\Big[e^{w_2s\frac{2\pi}3}+e^{w_4s\frac{2\pi}3}\Big],

$$

\begin{equation}

\label{mitr:eq22}

f_{3,2}(s)= \begin{vmatrix}

f_{3,2,11} (s) &f_{3,2,12} (s)\\

f_{3,2,21} (s) &f_{3,2,22}(s)

\end{vmatrix},

\end{equation}

%\allowdisplaybreaks

$$

f_{3,2,11} (s)=e^{w_1s\pi}+e^{w_3s\pi}-\alpha_1\Big[e^{w_1s\frac\pi3}+

e^{w_3s\frac\pi3}\Big]-\alpha_2

\Big[e^{w_1s\frac{2\pi}3}+e^{w_3s\frac{2\pi}3}\Big],$$

\begin{multline}

f_{3,2,12} (s)=A_{32}(\pi,s)+A_{34}(\pi,s)-\alpha_1\left[A_{32}\Bigl(\frac\pi3,s\Bigr)+

A_{34}\Bigl(\frac\pi3,s\Bigr)\right] -\\-\alpha_2 \left[A_{32}\Bigl(\frac{2\pi}3,s\Bigr)+A_{34}

\Bigl(\frac{2\pi}3,s\Bigr)\right],

\end{multline}

$$

f_{3,2,21} (s)=e^{w_1s\pi}+e^{w_3s\pi}-\beta_1\Big[e^{w_1s\frac\pi3}+

e^{w_3s\frac\pi3}\Big]-\beta_2\Big[e^{w_1s\frac{2\pi}3}+e^{w_3s\frac{2\pi}3}\Big],

$$

\begin{multline}

f_{3,2,22} (s)=A^2_{32}(\pi,s)+A^2_{34}(\pi,s)-\beta_1\left[A^2_{32}\Bigl(\frac\pi3,s\Bigr)+

A^2_{34}\left(\frac\pi3,s\right)\right] -\\-\beta_2 \left[A^2_{32}\Bigl(\frac{2\pi}3,s\Bigr)+A^2_{34}

\Bigl(\frac{2\pi}3,s\Bigr)\right],

\end{multline}

%\end{equation*}

функцию $f_6(s)$ можно вычислить аналогичным образом.





Основное приближение  уравнения \eqref{mitr:eq18}--\eqref{mitr:eq22} представляется в виде ${f_0(s){=}0}$ и от корней этого уравнения зависит поведение корней уравнения \eqref{mitr:eq18}--\eqref{mitr:eq22}. Для нахождения корней уравнения $f_0(s)=0$ необходимо изучить \textit{индикаторную диаграмму} этого уравнения (см.~\cite[гл.~12]{mitr:bib18}).



Введем обозначение $e^{\frac\pi3s}=z\ne0$, тогда основное приближение уравнения \eqref{mitr:eq18}--\eqref{mitr:eq22} примет следующий вид:

\begin{equation}

\label{mitr:eq23}

f_{0}(s)= \begin{vmatrix}

m_{11} (s) &m_{12} (s)\\

m_{21} (s) &m_{22}(s)

\end{vmatrix}=0,

\end{equation}

\begin{equation*}

\begin{gathered}

m_{11} (s)=z^{3w_1}-\alpha_2 z^{2w_1}-\alpha_1z^{w_1}-\alpha_1z^{-w_1}-\alpha_2z^{-2w_1}+z^{-3w_1},\\

m_{12} (s)=z^{3w_2}-\alpha_2 z^{2w_2}-\alpha_1z^{w_2}-\alpha_1z^{-w_2}-\alpha_2z^{-2w_2}+z^{-3w_2},\\

m_{21} (s)=z^{3w_1}-\beta_2 z^{2w_1}-\beta_1z^{w_1}-\beta_1z^{-w_1}-\beta_2z^{-2w_1}+z^{-3w_1},\\

m_{22} (s)=-z^{3w_2}+\beta_2 z^{2w_2}+\beta_1z^{w_2}+\beta_1z^{-w_2}+\beta_2z^{-2w_2}-z^{-3w_2}.

\end{gathered}

\end{equation*}



Раскладывая определитель $f_0(s)$ из \eqref{mitr:eq23} по столбцам на сумму определителей, получаем:

\begin{multline}

\label{mitr:eq24}

f_0(s)=(-2) z^{3w_1+3w_2}+\big(\alpha_2+\beta_2\big)z^{3w_1+2w_2}+

\big(\alpha_1+\beta_1\big) z^{3w_1+w_2}+ \\ +\big(\alpha_1+\beta_1\big) z^{3w_1-w_2}+

\big(\alpha_2+\beta_2\big)z^{3w_1-2w_2}+\big(-2\big) z^{3w_1-3w_2}+ \\ +

\big(\alpha_2+\beta_2\big)z^{2w_1+3w_2}-2\alpha_2\beta_2 z^{2w_1+2w_2}+\dots+

\big(\alpha_2+\beta_2\big) z^{2w_1-3w_2}+ \\ +\big(\alpha_1+\beta_1\big)z^{w_1+3w_2}-

\big(\alpha_1\beta_2+\alpha_2\beta_1\big)z^{w_1+2w_2}+\dots+   \\+\big(\alpha_1+\beta_1\big)z^{w_1-3w_2}+\dots+ \big(-2\big) z^{-3w_1+3w_2}+\big(\alpha_2+\beta_2\big)z^{-3w_1+2w_2}+ \\

+\big(\alpha_1+\beta_1\big)z^{-3w_2+w_2}+

\big(\alpha_1+\beta_1\big)z^{-3w_1-w_2}+  \\+\big(\alpha_2+\beta_2\big) z^{-3w_1-2w_2}+\big(-2\big)z^{-3w_1-3w_2}=0.

\end{multline}



Индикаторная диаграмма уравнения \eqref{mitr:eq23}, \eqref{mitr:eq24} "--- это выпуклая оболочка показателей экспонент, входящих в это уравнение. Учитывая формулы \eqref{mitr:eq4}, видим, что индикаторная диаграмма уравнения \eqref{mitr:eq23}, \eqref{mitr:eq24} имеет следующий вид:

\begin{equation}

\label{mitr:eq25}

\begin{array}{c}

\includegraphics{VestSamGTU} 

\end{array}.

\end{equation}

Индикаторная диаграмма представляет собой квадрат с вершинами в точках 

$$3w_1+3w_2,\quad  -3w_1+3w_2,\quad  -3w_1-3w_2\quad \text{и} \quad 3w_1-3w_2.$$ 

На стороне $\big[{3w_1-3w_2}; {3w_1+3w_2}\big]$ находятся точки (показатели экспонент) $3w_1-2w_2$, $3w_1-w_2$, $3w_1+w_2$, $3w_1+2w_2$, которые вместе с точкой $3w_1$ (не является показателем экспонент, входящих в \eqref{mitr:eq23}, \eqref{mitr:eq24}) делят этот отрезок на шесть равных частей. Аналогичная ситуация с другими сторонами квадрата. Если бы точки (показатели экспонент) делили отрезки на несоизмеримые части, то корни уравнения \eqref{mitr:eq23}, \eqref{mitr:eq24} и асимптотику корней уравнения \eqref{mitr:eq18}--\eqref{mitr:eq22} выписать было бы нельзя (см.~\cite{mitr:bib19,mitr:bib20}).



Корни уравнения \eqref{mitr:eq23}, \eqref{mitr:eq24} (а также корни уравнения \eqref{mitr:eq18}--\eqref{mitr:eq22}) находятся в четырех секторах бесконечно малого раствора, заштрихованных на рисунке \eqref{mitr:eq25}, биссектрисы которых являются срединными перпендикулярами к сторонам квадрата. Точки, находящиеся на сторонах квадрата ($3w_1-2w_2=3-2i$, $3w_1-w_2=3-i$, $3w_1+w_2=3+i,\dots$), также будут оказывать влияние  на поведение корней уравнений \eqref{mitr:eq23}, \eqref{mitr:eq24} и \eqref{mitr:eq18}--\eqref{mitr:eq22}. Точки, попадающие внутрь индикаторной диаграммы \eqref{mitr:eq25} ($2w_1-2w_2$, $2w_1-w_2$, $2w_1+w_2$, $2w_1+2w_2$, $w_1+2w_2$,$\dots$), на асимптотику корней уравнений \eqref{mitr:eq23}, \eqref{mitr:eq24} и \eqref{mitr:eq18}--\eqref{mitr:eq22} влиять не будут.



\smallskip



\Section{Асимптотика собственных значений в секторе 1) индикаторной диаграммы \eqref{mitr:eq25}}

Изучим сектор 1) индикаторной диаграммы \eqref{mitr:eq25}. Из общей теории (см. \cite[гл.~12]{mitr:bib18}, \cite{mitr:bib19, mitr:bib20}) следует, что в секторе 1) на асимптотику корней уравнений \eqref{mitr:eq23}, \eqref{mitr:eq24} и \eqref{mitr:eq18}--\eqref{mitr:eq22} влияют только точки, находящиеся на отрезке $[3w_1-3w_2;3w_1+3w_2]$, т.\,е. точки $3w_1-3w_2$, $3w_1-2w_2$, $3w_1-w_2$, $3w_1+w_2$, $3w_1+2w_2$ и $3w_1+3w_2$.



\smallskip



\begin{lemma}[4]Основное приближение уравнения на собственные значения дифференциального оператора \eqref{mitr:eq1}--\eqref{mitr:eq3} в секторе {\rm 1)} индикаторной диаграммы \eqref{mitr:eq25} имеет следующий вид\/{\rm:}

\begin{multline}

\label{mitr:eq26}

f_{0,1}(s)=(-2)z^{3w_1+3w_2}

+(\alpha_2+\beta_2)z^{3w_1+2w_2}

+(\alpha_1+\beta_1)z^{3w_1+w_2}+\\

+(\alpha_1+\beta_1)z^{3w_1-w_2}

+(\alpha_2+\beta_2)z^{3w_1-2w_2}

+(-2)z^{3w_1-3w_2}=0.

\end{multline}

\end{lemma}



\smallskip



В уравнении \eqref{mitr:eq26}, поделив  на $(-1) z^{3w_1}\ne0$, сделаем замену $z^{w_2}=x\ne0$. В итоге получим уравнение шестой степени:

\begin{equation*}

g_1(x)=2x^6-(\alpha_2+\beta_2)x^5

-(\alpha_1+\beta_1) x^4-

(\alpha_1+\beta_1)x^2-

(\alpha_2+\beta_2)x+2=0,

\end{equation*}

которое  в силу граничных условий \eqref{mitr:eq2} ($\alpha_1+\beta_1=-{10}/3$, $\alpha_2+\beta_2={16}/3$) принимает следующий вид:

\begin{equation}

\label{mitr:eq27}

g_1(x)=2x^6-\frac{16}3x^5

+\frac{10}3

 x^4+ \frac{10}x

x^2-\frac{16}3

x+2=0.

\end{equation}

Уравнение  \eqref{mitr:eq27} легко раскладывается на множители: 

$$g_1(x)=\frac23(x-1)^4(3x^2+4x+3)=0$$ и имеет следующие корни: 

\begin{equation}

\label{mitr:eq28}

\begin{gathered}

x_1=x_2=x_3=x_4=1 \ \ \text{(кратность равна четырем)};\\

x_5=-\frac{2}3+\frac{\sqrt 5}3i=e^{i\varphi}, \

x_6=\bar x_5=-\frac{2}3-\frac{\sqrt5}3i=e^{-i\varphi},  \\

\varphi=\arccos\Bigl(-\frac{2}3\Bigr)=\pi-\arcsin\Bigl(\frac{\sqrt5}3\Bigr).

\end{gathered}

\end{equation}



Формула \eqref{mitr:eq28} задает корни основного приближения уравнения \eqref{mitr:eq18}--\eqref{mitr:eq22} в~секторе 1) индикаторной диаграммы \eqref{mitr:eq25} (т.\,е. корни уравнений \eqref{mitr:eq26} и \eqref{mitr:eq27}). Из общей теории (см.~\cite[гл.~12]{mitr:bib18}) нахождения корней квазиполиномов вида \eqref{mitr:eq18}--\eqref{mitr:eq22} следует, что в секторе 1) надо оставить только те экспоненты, которые соответствуют правому вертикальному отрезку индикаторной диаграммы \eqref{mitr:eq25}. Поэтому для сектора 1) уравнение \eqref{mitr:eq17} принимает следующий вид:

\begin{gather*}

f_1(s)=

\begin{vmatrix}

b_{11} &b_{12}\\

b_{21} &b_{22}

\end{vmatrix}

=0,\\

b_{11}=z^{3w_1}-\frac{A_{31}(\pi,s)}{4s^3} +\underline O\Bigl(\frac1{s^6}\Bigr)+\bar{\bar o}(1),\\

z=e^{\frac\pi3s}\ne0, \\

b_{21}=w^2_1 z^{3w_1}-\frac{A^2_{31}(\pi,s)}{4s^3} +\underline O\Bigl(\frac1{s^6}\Bigr)+\bar{\bar o}(1), \\

w_1=1, \quad w^2_1=1, \\

b_{12}=\left[z^{3w_2}-\frac{A_{32}(\pi,s)}{4s^3} +\underline O\Bigl(\frac1{s^6}\Bigr)\right]-

\alpha_2\left[z^{2w_2}-\frac{A_{32}({2\pi}/3,s)}{4s^3} +\underline O\Bigl(\frac1{s^6}\Bigr)\right]- \\-

\alpha_1\left[z^{w_2}-\frac{A_{32}({\pi}/3,s)}{4s^3} +\underline O\Bigl(\frac1{s^6}\Bigr)\right]-

\alpha_1\left[z^{-w_2}-\frac{A_{34}({\pi}/3,s)}{4s^3} +\underline O\Bigl(\frac1{s^6}\Bigr)\right]- \\-

\alpha_2\left[z^{-2w_2}-\frac{A_{34}({2\pi}/3,s)}{4s^3} +\underline O\Bigl(\frac1{s^6}\Bigr)\right]+

\left[z^{-3w_2}-\frac{A_{34}({\pi},s)}{4s^3} +\underline O\Bigl(\frac1{s^6}\Bigr)\right],

\end{gather*}

\begin{equation}

\label{mitr:eq29}

w_2=i, \quad w_4=-w_2=-i, 

\end{equation}

\begin{gather*}

b_{22}=

\left[w^2_2 z^{3w_2}-\frac{A^2_{32}({\pi},s)}{4s^3}

+\underline O\Bigl(\frac1{s^6}\Bigr)\right]-

\beta_2\left[w^2_2z^{2w_2}-\frac{A^2_{32}({2\pi}/3,s)}{4s^3}

+\underline O\Bigl(\frac1{s^6}\Bigr)\right]- \\

-

\beta_1\left[w^2_2z^{w_2}-\frac{A^2_{32}({\pi}/3,s)}{4s^3}

+\underline O\Bigl(\frac1{s^6}\Bigr)\right]-

\beta_1\left[w^2_4z^{-w_2}-\frac{A^2_{34}({\pi}/3,s)}{4s^3}

+\underline O\Bigl(\frac1{s^6}\Bigr)\right]- \\

-

\beta_2\left[w^2_4z^{-2w_2}-\frac{A^2_{34}({2\pi}/3,s)}{4s^3}

+\underline O\Bigl(\frac1{s^6}\Bigr)\right]+

\left[w^2_4z^{-3w_2}-\frac{A^2_{34}({2\pi}/3,s)}{4s^3}

+\underline O\Bigl(\frac1{s^6}\Bigr)\right],  \\

w^2_2=w^2_4=-1. 

\end{gather*}



Изучая определитель $f_1(s)$ из \eqref{mitr:eq29}, сделаем необходимые преобразования и убедимся в справедливости следующего утверждения.



\smallskip



\begin{lemma}[5]Уравнение на собственные значения дифференциального оператора \eqref{mitr:eq1}--\eqref{mitr:eq3} в~секторе  {\rm  1)} индикаторной диаграммы  \eqref{mitr:eq25} имеет следующий вид\/{\rm:}

\begin{gather}

\label{mitr:eq30}

f_1(s)=f_{0,1}(s)-\frac{f_{3,1}(s)}{4s^3}+\frac{f_{6,1}(s)}{16s^6}+\underline O\Bigl(\frac1{s^9}\Bigr)=0,\\

\label{mitr:eq31}

f_{0,1}(s){=}

\begin{array}{|ll|}

z^{3w_1} &

z^{3w_2}-\alpha_2z^{2w_2}-\alpha_1z^{w_2}-\alpha_1z^{-w_2}-\alpha_2 z^{-2w_2}+z^{-3w_2}\\

z^{3w_1}

&{-}z^{3w_2}+\beta_2z^{2w_2}+\beta_1z^{w_2}+\beta_1z^{-w_2}+\beta_2 z^{-2w_2}-z^{-3w_2}

\end{array},\\

\label{mitr:eq34}

f_{3,1}=f_{3,1,1}(s)+f_{3,1,2}(s),\\

\label{mitr:eq34}

f_{3,1,1}(s){=}

\begin{array}{|ll|}

A_{31}(\pi,s)

&z^{3w_2}{-}\alpha_2z^{2w_2}{-}\alpha_1z^{w_2}{-}\alpha_1z^{-w_2}{-}\alpha_2 z^{-2w_2}{+}z^{-3w_2}\\

A^2_{31}(\pi,s)

&{-}z^{3w_2}{+}\beta_2z^{2w_2}{+}\beta_1z^{w_2}{+}\beta_1z^{-w_2}{+}\beta_2 z^{-2w_2}{-}z^{-3w_2}

\end{array},\\

\label{mitr:eq34}

f_{3,1,2}(s){=}

\begin{array}{|lr|}

z^{3w_1} &

A_{32}(\pi,s){-}\alpha_2A_{32}\big(\frac{2\pi}3,s\big){-}

\alpha_1A_{32}\big(\frac\pi3,s\big){-}

\alpha_1A_{34}\big(\frac\pi3,s\big){-}\\

&{-}\alpha_2A_{34}\big(\frac{2\pi}3,s\big)

{+}A_{34}(\pi,s)\\[2mm]

z^{3w_1} &

A^2_{32}(\pi,s){-}\beta_2A^2_{32}\big(\frac{2\pi}3,s\big){-}

\beta_1A^2_{32}\big(\frac\pi3,s\big){-}

\beta_1A^2_{34}\big(\frac\pi3,s\big)-\\

&{-}\beta_2A^2_{34}\big(\frac{2\pi}3,s\big)

{+}A^2_{34}(\pi,s)

\end{array}.

\end{gather}

\end{lemma}



\smallskip



Действительно, раскладывая определитель $f_1(s)$ из \eqref{mitr:eq30} по столбцам на сумму определителей и упорядочивая величины по росту $s$, получаем \eqref{mitr:eq30}--\eqref{mitr:eq34}, величины $A_{31}(\pi,s)$, $A^2_{31}(\pi,s)$, $A_{32}(\pi,s),\dots$, $A^2_{34}(\pi,s)$ определены формулами \eqref{mitr:eq12}, \eqref{mitr:eq13}.



Основное приближение уравнения \eqref{mitr:eq30}--\eqref{mitr:eq34} имеет вид $f_{0,1}(s)=0$, где $f_{0,1}(s)$ определено формулой \eqref{mitr:eq31}, и это уравнение преобразуется к виду 

\begin{multline}

f_{0,1}(s)=z^{3w_1}\Big[(-2)z^{3w_2}+\big(\alpha_2+\beta_2\big)z^{2w_2}+

\big(\alpha_1+\beta_1\big)z^{w_2} +

\big(\alpha_1+\beta_1\big)z^{-w_2}+ \\+

\big(\alpha_2+\beta_2\big)z^{-2w_2}

+(-2)z^{-3w_2}\Big]=0,\quad  z^{3w_1}\ne0.

\end{multline}

Это уравнение при выполнении условия \eqref{mitr:eq2} принимает следующий вид:

\begin{equation*}

\begin{gathered}

3z^{3w_2}-8z^{2w_2}+ 5z^{w_2}+5z^{-w_2}-

8z^{-2w_2}+3z^{-3w_2}=0, \\

3x^6-8x^5+5x^4+5x^2-8x+3=0, \quad  x=z^{w_2}=

e^{a\frac\pi3 w_2s}\ne0,

\end{gathered}

\end{equation*}

т.\,е. совпадает с уравнением $g_1(x)=0$ из \eqref{mitr:eq27}, которое имеет корни, описываемые формулой \eqref{mitr:eq28}.





Изучим сначала наиболее интересный случай:

\begin{equation}

\label{mitr:eq35}

\begin{gathered}

x_1=x_2=x_3=x_4=1 \quad \text{(кратность 4)} \Leftrightarrow

e^{\frac\pi3 w_2s}=1=e^{2\pi ik}, \\ 

k\in\mathbb{N}

\Leftrightarrow s_{k,1,\text{осн}}=\frac{6ki}{w_2}\stackrel{\eqref{mitr:eq4}}{=}

{6k}

 \ \text{(кратность 4),} \ k\in\mathbb{N}.

%\eqno\eqref{mitr:eq45}

\end{gathered}

\end{equation}



В формуле \eqref{mitr:eq35} у серии корней $s_{k,1,\text{осн}}$ индекс~1 показывает, что мы рассматриваем сектор 1) индикаторной диаграммы \eqref{mitr:eq25}, <<осн>> означает <<основное приближение>>. Из общей теории нахождения корней квазиполиномов вида \eqref{mitr:eq30}--\eqref{mitr:eq34} в случае условия \eqref{mitr:eq28} (см.~\cite[гл.~12]{mitr:bib18}, \cite{mitr:bib11, mitr:bib13, mitr:bib20}) следует справедливость следующего утверждения.



\smallskip



\hypertarget{mitr:th1}{}



\begin{theorem}[1]Асимптотика собственных значений дифференциального оператора \eqref{mitr:eq1}--\eqref{mitr:eq3} в~секторе~{\rm1)} индикаторной диаграммы  \eqref{mitr:eq25} имеет вид

\begin{equation}

\label{mitr:eq36}

s_{k,1,p}={6k}+\frac{d_{1kp}}{k^{3/4}}+

\frac{d_{2kp}}{k^{6/4}}+

\frac{d_{3kp}}{k^{9/4}}+

\frac{d_{4kp}}{k^{12/4}}+

\underline O\Bigl(\frac1{k^{15/4}}\Bigr), \quad

p=1,2,3,4.

\end{equation}

\end{theorem}



\smallskip



\begin{proof} Для доказательства теоремы~\hyperlink{mitr:th1}{1}  необходимо показать, что коэффициенты $d_{1kp}$, $d_{2kp}$, $d_{3kp}$, $\dots$  $(p=1,2,3,4)$ формулы \eqref{mitr:eq36} находятся в явном виде единственным образом.



Применяя формулы Маклорена, имеем

%\begin{equation}

\begin{multline*}

z^{\pm w_2}

\big|_{s_{k,1,p}}

=e^{\pm s\frac\pi3 w_2}

\big|_{s_{k,1,p}}=1\pm \frac\pi3 iM_{kp}-\frac{\pi^2}9 M^2_{kp}\mp\frac{\pi^3 i}{27} M^3_{kp}+\\

+\frac{\pi^4}{81} M^4_{kp} \pm\frac{\pi^5 i}{243} M^5_{kp}-\dots,

\end{multline*}

\begin{equation}

\label{mitr:eq37}

M_{kp}=\frac{d_{1kp}}{k^{3/4}} +\frac{d_{2kp}}{k^{6/4}}+\frac{d_{3kp}}{k^{9/4}}+\frac{d_{4kp}}{k^{12/4}}+\underline O\Bigl(\frac1{k^{{15}/4}}\Bigr),

\end{equation}

\begin{multline}

\label{mitr:eq39}

z^{\pm 2w_2}

\big|_{s_{k,1,p}}=

1\pm \frac{2\pi}3 iM_{kp}-\frac{4\pi^2}9 M^2_{kp}\mp

\frac{8\pi^3}{27}i M^3_{kp}+\\

+\frac{16\pi^4}{81} M^4_{kp}\pm

\frac{32\pi^5}{243} iM^3_{kp}-\dots,

\end{multline}

\begin{equation}

\label{mitr:eq39}

z^{\pm 3w_2}

\big|_{s_{k,1,p}}=

1 \pm \pi i M_{kp}-\pi^2 M^2_{kp}\mp \pi^3 i M^3_{kp}+\pi^4M^4_{kp}\pm \pi^5 iM^5_{kp}-\dots,

\end{equation}

\begin{equation}

\label{mitr:eq40}

\begin{gathered}

\frac1{s^3}

\Big|_{s_{k,1,p}}=\frac{1}{6^3k^3}\left(1-\frac{3d_{1kp}}{6k^{7/4}}+\underline O\Bigl(\frac1{k^{10/4}}\Bigr)\right),  \\

\frac1{s^6_{k,1,p}} =\frac{1}{6^6k^6}\left(1-\frac{6 d_{1kp}}{6k^{7/4}}

+\underline O\Bigl(\frac1{k^{10/4}}\Bigr)\right).

\end{gathered}

\end{equation}





Подставляя формулы \eqref{mitr:eq37}--\eqref{mitr:eq40} и \eqref{mitr:eq36} в уравнение \eqref{mitr:eq30}--\eqref{mitr:eq34}, видим, что коэффициенты при $M^0_{kp}$, $M^1_{kp}$, $M^2_{kp}$, $M^3_{kp}$ равны нулю в силу уравнения \eqref{mitr:eq27}, коэффициент при $M^4_{kp}$ равен $-\frac{160\pi^4}{81}$. В итоге мы получаем

\begin{multline}

\label{mitr:eq41}

-\frac{160\pi^4}{81}M^4_{kp}+0\cdot M^5_{kp}+\gamma_6M^6_{kp}+\dots-\\-

\frac1{24k^3}\Bigl(1-\frac{3d_{1kp}}{6k^{7/4}}

+\underline O\Bigl(\frac1{k^{10/4}}\Bigr)\Bigr) G_3(s)\Big|_{s_{k,1,k,p}} +

\underline O\Bigl(\frac1{k^6}\Bigr)=0, \quad \gamma_6\ne0,

\end{multline}

\begin{multline}

\label{mitr:eq42}

G_3(s)=\int ^\pi_0 q(t)dt_{a11}

\Big[-2\cdot z^{3w_2}+ \big(\alpha_2+\beta_2\big) z^{2w_2}+\big(\alpha_1+\beta_1\big)z^{w_2}+\\+

\big(\alpha_1+\beta_1\big)z^{-w_2}+

\big(\alpha_2+\beta_2\big)z^{-2w_2}-2\cdot z^{-3w_2}\Big]+G_2(s),

\end{multline}

%\allowdisplaybreaks

%\begin{equation}

\begin{multline}

G_2(s) =(-2)w_2z^{3w_2}\varphi_{22}(\pi,s)-

2w_4z^{-3w_2}\varphi_{24}(\pi,s)+ \\

+\big(\alpha_2+\beta_2\big) w_2z^{2w_2}\varphi_{22}\Bigl(\frac{2\pi}3,s\Bigr)+

\big(\alpha_2+\beta_2\big) w_4z^{-2w_2}\varphi_{24}\Bigl(\frac{2\pi}3,s\Bigr)+\\ +

\big(\alpha_1+\beta_1\big) w_2z^{w_2}\varphi_{22}\Bigl(\frac{\pi}3,s\Bigr)+

\big(\alpha_1+\beta_1\big) w_4z^{-w_2}\varphi_{24}\Bigl(\frac{\pi}3,s\Bigr)+\\ +

\big(\alpha_1+\beta_1\big) w_2z^{w_2}\varphi_{42}\Bigl(\frac{\pi}3,s\Bigr)+

\big(\alpha_1+\beta_1\big) w_4z^{-w_2}\varphi_{44}\Bigl(\frac{\pi}3,s\Bigr)+\\+

\big(\alpha_2+\beta_2\big) w_2z^{2w_2}\varphi_{42}\Bigl(\frac{2\pi}3,s\Bigr)+

\big(\alpha_2+\beta_2\big) w_4z^{-2w_2}\varphi_{44}\Bigl(\frac{2\pi}3,s\Bigr)-\\-

2w_2z^{3w_2}\varphi_{42}(\pi,s)- 2w_4z^{-3w_2}\varphi_{44}(\pi,s),

\label{mitr:eq43}

\end{multline}

где

\begin{equation}

w_2\stackrel{\eqref{mitr:eq4}}{=}i,  \quad w_4\stackrel{\eqref{mitr:eq4}}{=}-i; \quad \alpha_1+\beta_1 \stackrel{\eqref{mitr:eq2}}{=}-\frac{10}3, \quad \alpha_2+\beta_2 \stackrel{\eqref{mitr:eq2}}{=}\frac{16}3,

\end{equation}

а функции $\varphi_{kn}(x,s)$ определены формулой \eqref{mitr:eq10}.





Подставляя  формулы \eqref{mitr:eq37}--\eqref{mitr:eq39} в формулу \eqref{mitr:eq43}, имеем

\begin{multline}

\label{mitr:eq44}

G_2(s)\big|_{s_{k,1,p}}=

(-2)w_2\Big[1+\pi i M_{kp}-\pi^2 M^2_{kp}+\underline O(M^3_{kp})\Big]

\varphi_{22}({\pi},s_{s_{k,1,p}})+\\+

2w_2\Big[1-\pi i M_{kp}-\pi^2 M^2_{kp}+\underline O(M^3_{kp})\Big]

\varphi_{44}({\pi},s_{s_{k,1,p}})+\\

+

\big(\alpha_2+\beta_2\big)w_2

\Bigl[1+\frac{2\pi}3 i M_{kp}-\frac{4\pi^2}9 M^2_{kp}+\underline O(M^3_{kp})\Bigr]

\varphi_{22}\Bigl(\frac{2\pi}3,s_{s_{k,1,p}}\Bigr)-\\-

\big(\alpha_2+\beta_2\big)

w_2\Bigl[1-\frac{2\pi}3 i M_{kp}-\frac{4\pi^2}9 M^2_{kp}+\underline O(M^3_{kp})\Bigr]

\varphi_{44}\Bigl(\frac{2\pi}3,s_{s_{k,1,p}}\Bigr)+\\

+\dots+H_2(s)\big|_{s_{k,1,p}}, \quad

p=1,2,3,4,

\end{multline}

\begin{multline}

\label{mitr:eq45}

H_2(s)\big|_{s_{k,1,p}}=

2w_2\Big[z^{-3w_2}\varphi_{24}\big(\pi,s_{s_{k,1,p}}\big)- z^{3w_2}

\varphi_{42}\big(\pi,s_{s_{k,1,p}}\big)\Big]-\\-

\big(\alpha_2+\beta_2\big)

w_2\Bigl[z^{-w_2}\varphi_{24}\Bigl(\frac{2\pi}3,s_{s_{k,1,p}}\Bigr)- z^{2w_2}

\varphi_{42}\Bigl(\frac{2\pi}3,s_{s_{k,1,p}}\Bigr)\Bigr]+\\

+

\big(\alpha_1+\beta_1\big)

w_2\Bigl[z^{-w_2}\varphi_{24}\Bigl(\frac{\pi}3,s_{s_{k,1,p}}\Bigr)- z^{w_2}

\varphi_{42}\Bigl(\frac{2\pi}3,s_{s_{k,1,p}}\Bigr)\Bigr].

\end{multline}

Проведя в формуле \eqref{mitr:eq44} необходимые вычисления и преобразования и~учитывая, что  $$\varphi_{22}(\pi,s)\stackrel{\eqref{mitr:eq10}}{=} \varphi_{42}(\pi,s)\stackrel{\eqref{mitr:eq10}}{=}\int^\pi_0 q(t) dt,$$ получим

\begin{equation}

\label{mitr:eq46}

G_2(s)\big|_{s_{k,1,p}}=\frac89 w_2 \pi i\int^\pi_0 q(t)dt

\Bigl[\frac{d_{1kp}}{k^{3/4}}+\frac{d_{2kp}}{k^{6/4}}+\underline O

\Bigl(\frac1{k^{9/4}}\Bigr)+H_2(s)\Bigr]\Big|_{s_{k,1,p}}.

\end{equation}



Учитывая, что $z=e^{\frac \pi 3 s}$, $w_2=i$, $w_4=-i$, и подставляя $s_{k,1,p}$ из \eqref{mitr:eq36} в~формулу \eqref{mitr:eq45}, имеем

%\allowdisplaybreaks

\begin{multline}

\label{mitr:eq47}

\big[z^{-3w_2}\varphi_{24}(\pi,s)- z^{3w_2}

\varphi_{42}(\pi,s)\big]\big |_{s_{k,1,p}}

\stackrel{\eqref{mitr:eq10}}{=}  \\ 

=

\exp\Bigl[\frac{\pi i}3 (-3)\big({6k}+M_{kp}\big)\Bigr]

\int ^\pi_0 q(t)\exp

\Big[\big(w_2-w_4\big)\big({6k}+M_{kp}t\big)\Big]dt-\\-

\exp\Bigl[\frac{\pi i}3 3\big({6k}+M_{kp}\big)\Bigr]

\int ^\pi_0 q(t)\exp

\Big[-\big(w_2-w_4\big)\big({6k}+M_{kp}\big)t\Big]dt=\\ 

=2i\int ^\pi_0 q(t)\sin\big[12kt+(2t-\pi)M_{kp}\big]dt.

\end{multline}

Аналогичным образом выводятся формулы

\begin{multline}

\label{mitr:eq48}

\Bigl[z^{-2w_2}\varphi_{24}\Bigl(\frac{2\pi}3,s\Bigr)- z^{2w_2}\varphi_{42}\Bigl(\frac{2\pi}3,s\Bigr)\Bigr]\Big|_{s_{k,1,p}}=\\

=

2i\int ^{2\pi/3}_0q(t)\sin\left[12kt+\Bigl(2t-\frac{2\pi}3\Bigr)M_{kp}\right]dt;

\end{multline}

\begin{multline}

\label{mitr:eq48}

\Bigl[z^{-w_2}\varphi_{24}\Bigl(\frac{\pi}3,s\Bigr)-

z^{w_2}\varphi_{42}\Bigl(\frac{\pi}3,s\Bigr)\Bigr]\Big|_{s_{k,1,p}}= \\=

2i\int ^{\pi/3}_0q(t)\sin\left[12kt+\Bigl(2t-\frac\pi3\Bigr)M_{kp}\right]dt.

\end{multline}

Подставляя формулы \eqref{mitr:eq47}--\eqref{mitr:eq48} в \eqref{mitr:eq45}, получаем

\begin{multline}

\label{mitr:eq49}

H_2(s)\big|_{s_{k,1,p}}=(-4)\int ^\pi_0 q(t)\sin\big[12kt+(2t-\pi)M_{kp}\big]dt+\\ +

\frac{32}3\int ^{2\pi/3}_0q(t)\sin\left[12kt+\Bigl(2t-\frac{2\pi}3\Bigr)M_{kp}\right]dt+\\+

\frac{20}3\int ^{\pi/3}_0q(t) \sin\left[12kt+\Bigl(2t-\frac\pi3\Bigr)M_{kp}\right]dt .

\end{multline}

Подставляя формулы \eqref{mitr:eq46}, \eqref{mitr:eq49} в \eqref{mitr:eq41}, \eqref{mitr:eq42}, находим

\begin{multline}

\label{mitr:eq50}

\frac{d^4_{1kp}}{k^{12/4}} +

\frac{4d^3_{1kp}d_{2kp}}{k^{15/4}}+

\underline O\Bigl(\frac1{k^{18/4}}\Bigr)=\\=

\frac{3^4}{15\cdot 2^8}\frac1{k^{12/4}}

\biggl\{

\pi i w_2\frac{d_{1kp}}{k^{3/4}}\biggl[

(-4)\int ^\pi_0 q(t)dt+

\frac{32}3\int ^{2\pi/3}_0q(t)dt + \\+

\frac{20}3\int ^{\pi/3}_0 q(t)dt\biggr]+ 

\left[\frac{20}3v_1\Bigl(\frac\pi3,k\Bigr)+  \frac{32}3v_1\Bigl(\frac{2\pi}3,k\Bigr)-

4v_1(\pi,k)\right]+\\+

\frac{d_{1kp}}{k^{3/4}}

\left[\frac{20}3v_2\Bigl(\frac\pi3,k\Bigr)+

\frac{32}3v_3\Bigl(\frac{2\pi}3,k\Bigr)-

4v_4(\pi,k)\right]

+\underline O\Bigl(\frac1{k^{6/4}}\Bigr)\biggr\},

\end{multline}

где введены следующие обозначения:

\begin{equation}

\label{mitr:eq51}

\begin{gathered}

v_1(b,k)=\int^b_0 q(t) \sin(12kt)dt; \hspace{2cm} \\

v_2(b,k)=\int^b_0 q(t)\Bigl(2t-\frac\pi3\Bigr) \cos(12kt)dt; \; \; \\

v_3(b,k)=\int^b_0 q(t)\Bigl(2t-\frac{2\pi}3\Bigr) \cos(12kt)dt;  \\

v_4(b,k)=\int^b_0 q(t)\Bigl(2t-\frac{3\pi}3\Bigr) \cos(12kt)dt.

%\eqno\eqref{mitr:eq61}

\end{gathered}

\end{equation}



Приравнивая в формуле \eqref{mitr:eq50} коэффициенты при $k^{-{12}/4}$, получим

\begin{equation}

\label{mitr:eq52}

d_{1kp}=\frac{3w_p}{2\sqrt[4]{15}}

\sqrt[4]{F_k(q)}, \quad  k\in \mathbb{N}, \quad

p=1,2,3,4, \quad w_p=e^{\frac{2\pi i}4(p-1)},

\end{equation}

\begin{multline}

F_k(q)=\frac{20}3\int ^{\pi/3}_0 q(t)\sin (12kt)dt+ \\+

\frac{32}3\int ^{2\pi/3}_0 q(t)\sin (12kt)dt-4\int ^\pi_0 q(t)\sin

(12kt)dt.

\end{multline}



Таким образом, мы видим, что кратные в главном приближении корни $x_1=x_2=x_3=x_4=1$ (кратности 4) (и $s_{k,1,\text{осн}}\stackrel{\eqref{mitr:eq35}}{=}{6k}$ кратности 4) <<расщепляются>> на четыре серии однократных собственных значений вида~\eqref{mitr:eq36},~\eqref{mitr:eq52}.



Приравнивая в формуле \eqref{mitr:eq50} коэффициенты при  $k^{-{15}/4}$ и проводя необходимые преобразования и упрощения, выводим

\begin{multline}

\label{mitr:eq53}

d_{2kp}=

-\frac{27}{20\cdot 2^8 }\frac1{d^2_{1kp}}

\biggl\{\biggl[

\frac{20}3\int ^{\pi/3}_0 q(t)dt+

\frac{32}3\int ^{2\pi/3}_0q(t)dt-4\int ^\pi_0q(t) dt\biggr]+\\+

\left[\frac{40}3v_5\Bigl(\frac\pi3,k\Bigr)+ \frac{64}3v_5\Bigl(\frac{2\pi}3,k\Bigr) -8v_5(\pi,k)\right]-\\ -\frac\pi3\left[ \frac{20}3v_6\Bigl(\frac\pi3,k\Bigr)+ \frac{64}3v_6\Bigl(\frac{2\pi}3,k\Bigr) -12v_6(\pi,k)\right]\biggr\}, \\

k\in \mathbb{N}, \quad p=1,2,3,4,

\end{multline}

где

\begin{gather*}

v_5(b,k)=\int^b_0 t q(t)\cos(12kt)dt,\quad 

v_6(b,k)=\int^b_0  q(t)\cos(12kt)dt,

\end{gather*}

а величины $d_{1kp}$ определены формулой \eqref{mitr:eq52}.



Получение формул \eqref{mitr:eq52}, \eqref{mitr:eq53} завершает доказательство теоремы~\hyperlink{mitr:th1}{1}.

\end{proof}



\smallskip



Нахождение асимптотики собственных значений, соответствующих корням $x_{5,6}=e^{\pm i\varphi}$ формулы \eqref{mitr:eq28} (кратности 1), происходит согласно методике, продемонстрированной автором в работах~\cite{mitr:bib4,mitr:bib21,mitr:bib22}. Справедлива следующая теорема.



\smallskip



\hypertarget{mitr:th2}{}



\begin{theorem}[2]Асимптотика собственных значений дифференциального оператора \eqref{mitr:eq1}--\eqref{mitr:eq3} {\rm(}в секторе~{\rm1)} индикаторной диаграммы {\rm \eqref{mitr:eq25},} соответствующих однократным корням $x_{5,6}=e^{\pm i\varphi}$ уравнения {\rm \eqref{mitr:eq28},} имеет следующий вид\/{\rm:}

\begin{equation}

\label{mitr:eq54}

s_{k,1,m}={6\tilde k}+\frac{d_{3km}}{\tilde k^3}

+\frac{d_{6km}}{\tilde k^6}+

\underline O\Bigl(\frac1{\tilde k^9}\Bigr), \quad \tilde k=k\pm\frac{\varphi}{2\pi}, \quad m=5,6.

\end{equation}

\end{theorem}



\smallskip



\begin{proof} Действительно, уравнение на собственные значения оператора \eqref{mitr:eq1}--\eqref{mitr:eq3}  в секторе 1) имеет вид \eqref{mitr:eq30}--\eqref{mitr:eq34}, основное приближение которого (уравнение \eqref{mitr:eq27}) имеет  корни, заданные формулой \eqref{mitr:eq28}: $x_1=x_2\hm=x_3=x_4=1$

 (кратность 4, случай уже разобран в \eqref{mitr:eq36}, \eqref{mitr:eq52}, \eqref{mitr:eq53}) и~$x_5=e^{i\varphi}$, $x_6=e^{-i\varphi}$, $\varphi=\arccos (-\frac{2}3)=\pi-\arcsin (\frac{\sqrt5}3)\in(\frac\pi2;\pi)$, кратность равна~1, $x\stackrel{\eqref{mitr:eq26}}{=}z^{w_2}=e^{\frac\pi3 is}\ne0$.



 Поэтому 

 $$x_{5,6}=e^{\pm i\varphi} \Leftrightarrow e^{\frac\pi3 w_2s}=

e^{\pm i\varphi}e^{2\pi ik},\quad k\in\mathbb{N},$$ 

$$w_2=i \Leftrightarrow s_{k,1,m}={6k}\pm\frac{3\varphi}{\pi}=6 \tilde k,\quad \tilde k=k\pm\frac\varphi{2\pi},\quad  k\in\mathbb{N},\quad m=5,6.$$



Значит,  из общей теории нахождения корней квазимногочленов вида \eqref{mitr:eq30}--\eqref{mitr:eq34} (см.~\cite{mitr:bib11,mitr:bib13,mitr:bib20}) следует, что асимптотику собственных значений необходимо искать в~виде \eqref{mitr:eq54}. Для нахождения коэффициентов $d_{3km}$ $(m=5,6)$ в~явном виде действуем аналогично выкладкам \eqref{mitr:eq37}--\eqref{mitr:eq53}. Используя формулы Тейлора, имеем

\begin{multline}

\label{mitr:eq55}

z^{\pm nw_2}\big|_{s_{k,1,m}}=

e^{\pm n\frac\pi3 w_2s} \big|_{s_{k,1,m}} =\\=

\exp\Bigl(\pm n\frac\pi3 w_2{6\tilde k}\Bigr)

\exp\Bigl[\pm n\frac\pi3 w_2\Bigl(\frac{d_{3km}}{\tilde k^3}+\underline O

\Bigl(\frac1{\tilde k^6}\Bigr)\Bigr)\Bigr]=\\=

e^{\pm ni\varphi}\Bigl[1\pm \frac{n\pi w_2}3\frac{d_{3km}}{\tilde k^3}+

\underline O\Bigl(\frac1{\tilde k^6}\Bigr)\Bigr], \quad n=1,2,3.

\end{multline}



Подставляя формулы \eqref{mitr:eq54}, \eqref{mitr:eq55} в уравнение \eqref{mitr:eq30}--\eqref{mitr:eq34}, видим, что коэффициент при $\tilde k^0$ равен нулю в силу того, что $x_{5,6}=e^{\pm i\varphi}$ являются корнями уравнений \eqref{mitr:eq27} и $f_{0,1}=0$ из \eqref{mitr:eq30}, \eqref{mitr:eq31}. Приравнивая коэффициенты при $\tilde k^{-3}$, находим

\begin{equation}

\label{mitr:eq56}

d_{3km}=\frac{81\sqrt5}{20\pi \cdot 6^3\cdot 424}\frac{G_3(s)}{z^{3w_1}}\Big|_{s_{k,1,m,\text{осн}}}, \quad  m=5, \quad m=6,

\end{equation}

где

функция $G_3(s)$ определена в \eqref{mitr:eq41}--\eqref{mitr:eq43}. 

Произведя необходимые преобразования в \eqref{mitr:eq56}, получаем

\begin{multline}

\label{mitr:eq57}

d_{3km}=\frac{3\sqrt5}{160\cdot 424\pi}\biggl[

4\sin(3\varphi)\int ^\pi_0 q(t)dt-

\frac{32}3\sin (2\varphi)\int ^{2\pi/3}_0q(t)dt+\\+

\frac{20}3\sin \varphi\int ^{\pi/3}_0 q(t)dt-

4\int ^\pi_0 q(t)\sin\big(12\tilde kt-3\varphi\big)dt+\\+

\frac{32}3\int ^{2\pi/3}_0q(t)

\sin\big(12\tilde kt-2\varphi\big)dt-\\

-\frac{20}3\int ^{\pi/3}_0q(t)\sin\big(12\tilde kt-\varphi\big)dt\biggr], \quad m=5, \quad m=6.

\end{multline}



Если в уравнении \eqref{mitr:eq30}--\eqref{mitr:eq34} вести вычисления с точностью формул \eqref{mitr:eq12}, \eqref{mitr:eq13} (а не только \eqref{mitr:eq10}, \eqref{mitr:eq11}), то можно в явном виде вычислить коэффициенты $d_{6km}$ формулы \eqref{mitr:eq54}. Получение формул \eqref{mitr:eq57} завершает доказательство теоремы~\hyperlink{mitr:th2}{2} (доказано, что коэффициенты $d_{3km}$ формулы \eqref{mitr:eq54} находятся единственным образом и приведены явные формулы для их вычисления).

\end{proof}



\smallskip



Изучая аналогичным образом сектора 2), 3), 4) индикаторной диаграммы \eqref{mitr:eq25}, доказываем следующее утверждение.



\smallskip



\hypertarget{mitr:th3}{}



\begin{theorem}[3]Асимптотика собственных значений дифференциального оператора \eqref{mitr:eq1}--\eqref{mitr:eq3} в~секторах~{\rm2), 3), 4)} индикаторной диаграммы \eqref{mitr:eq25} подчиняется следующим законам\/{\rm:}

\begin{equation}

\label{mitr:eq59}

1)\hskip1.5cm \hfill s_{k,n,p}=s_{k,1,p}e^{\frac{\pi i}2(n-1)}, \ n=2,3,4, \ p=1,2,3,4; \

\lambda_{k,n,p}=s^4_{k,n,p}, \hfill

\end{equation}



~~$s_{k,1,p}$ определены формулами \eqref{mitr:eq36}, \eqref{mitr:eq52}, \eqref{mitr:eq53};

\begin{equation}

\label{mitr:eq59}

2)\quad s_{k,n,m}=s_{k,1,m}e^{\frac{2\pi i}4(n-1)}, \ n=2,3,4, \ m=5, \ m=6; \

\lambda_{k,n,m}=s^4_{k,n,m},

\end{equation}



~~$s_{k,1,m}$ определены формулами \eqref{mitr:eq54}, \eqref{mitr:eq57}.

\end{theorem}



Формулы \eqref{mitr:eq36}, \eqref{mitr:eq52}--\eqref{mitr:eq54}, \eqref{mitr:eq57}--\eqref{mitr:eq59} позволяют вычислить асимптотику собственных функций многоточечного дифференциального оператора \eqref{mitr:eq1}--\eqref{mitr:eq3}.



Аналогичное поведение спектра (с эффектом <<расщепления>> кратных в~главном приближении собственных значений) наблюдается у дифференциального оператора, задаваемого дифференциальным уравнением \eqref{mitr:eq1} с граничными условиями

\begin{gather}

y'(0)=0; \quad y^{(3)}(0)=0; \quad y(\pi)=\alpha_1 y\Bigl(\frac\pi3\Bigr)+\alpha_2 y\Bigl(\frac{2\pi}3\Bigr);\notag \\

\label{mitr:eq60}

 y''(\pi)=\beta_1y''\left(\frac\pi3\right)+\beta_2y''\Bigl(\frac{2\pi}3\Bigr); \\

\beta_1=2-\alpha_1, \quad \alpha_1\in\mathbb{C}; \quad \beta_2=-\alpha_2, \quad \alpha_2\in\mathbb{C},\notag

\end{gather}

с условием \eqref{mitr:eq3} суммируемости потенциала $q(x)$.



В этом случае уравнение \eqref{mitr:eq26}, \eqref{mitr:eq27} (основное приближение уравнения \eqref{mitr:eq18}--\eqref{mitr:eq22}) принимает следующий вид:

\begin{equation}

\label{mitr:eq61}

g_2(x)=2x^6-2x^4-2x^2+2=0, \quad

x=z^{w_2}=e^{\frac\pi3 is}\ne0.

\end{equation}



Уравнение \eqref{mitr:eq61} имеет следующие корни ($g_2(x)=2(x-1)^2(x+1)^2(x^2+1)$): $x_1=x_2=1$ (кратность 2), $x_3=x_4=-1$ (кратность 2), $x_5=i$ (кратность~1), $x_6=-i$ (кратность~1). Поэтому корни уравнения \eqref{mitr:eq18}--\eqref{mitr:eq22} (в основном приближении) имеют следующий вид:

\begin{gather*}

x_1=x_3=1 \Leftrightarrow e^{\frac\pi3is}=1=e^{2\pi ik}\Leftrightarrow s_{k,j,\text{осн}} =6k \ \ \text{кратность 2},\\ \ k\in\mathbb{N}, \ j=1,2;\\

x_3=x_4=-1 \Leftrightarrow e^{\frac\pi3is}=-1=e^{\pi i}e^{2\pi ik}\Leftrightarrow s_{k,j,\text{осн}} =6k +3\ \ \text{кратность 2},\\ \ k\in\mathbb{N}, \ j=3,4;\\

x_{5,6}=\pm i \Leftrightarrow e^{\frac\pi3is}=e^{\pm i\frac\pi2}e^{2\pi ik}\Leftrightarrow s_{k,j,\text{осн}} =6k\pm \frac32 \ \ \text{кратность 1},\\ \ k\in\mathbb{N}, \ j=5,6.

\end{gather*}



Справедливы утверждения, аналогичные теоремам~\hyperlink{mitr:th1}{1}, \hyperlink{mitr:th2}{2}, \hyperlink{mitr:th3}{3}.



\smallskip



\begin{theorem}[4] Асимптотика собственных значений дифференциального оператора  {\rm \eqref{mitr:eq1}, \eqref{mitr:eq60}, \eqref{mitr:eq3}} в секторах  {\rm  1), 2), 3), 4)} индикаторной диаграммы \eqref{mitr:eq25} имеет следующий вид\/{\rm:}

\begin{gather}

\label{mitr:eq62}

\hskip-2.5cm 1)\hskip2.5cm \hfill s_{k,1,j}=6k+\frac{d_{1kj}}{k^{3/2}}+\frac{d_{2kj}}{k^{6/2}}+\frac{d_{3kj}}{k^{9/2}}+

\underline O\Bigl(\frac1{k^{12/2}}\Bigr),\hfill \\

 \ k\in\mathbb{N}, \ j=1,2; \  s_{k,n,j}=s_{k,1,j} e^{\frac{\pi i}2(n-1)}, \ n=1,2,3,4; \notag\\ 

\lambda_{k,n,j}=s^4_{k,n,j}, \ j=1,2; \ n=1,2,3,4;\notag

\end{gather}

\begin{gather}

\hskip-2.5cm 2)\hskip1cm s_{k,1,j}=6k+3+\frac{d_{1kj}}{(6k+3)^{3/2}}+\frac{d_{2kj}}{(6k+3)^{6/2}}+\hskip1.5cm \hfill \notag\\

\label{mitr:eq64}

\hskip3cm 

+\frac{d_{3kj}}{(6k+3)^{9/2}}+

\underline O\Bigl(\frac1{(6k+3)^{12/2}}\Bigr),\\ 

k\in\mathbb{N}, \quad j=3,4; \quad  s_{k,n,j}=s_{k,1,j} e^{\frac{\pi i}2(n-1)}, \quad n=1,2,3,4; \notag\\ 

\lambda_{k,n,j}=s^4_{k,n,j}, \quad j=3,4; \quad n=1,2,3,4;\notag

\end{gather}

\begin{gather}

\label{mitr:eq64}

\hskip-1.8cm 3)\hskip2.5cm

s_{k,1,j}=6\tilde {\tilde k}+\frac{d_{1kj}}{\tilde {\tilde k}^{3}}+\frac{d_{2kj}}{\tilde {\tilde k}^{6}}+

\underline O\Bigl(\frac1{\tilde {\tilde k}^{12/2}}\Bigr), \ \tilde {\tilde k}=k+\frac14,\hfill \\ j=5,6; \  s_{k,n,j}=s_{k,1,j} e^{\frac{\pi i}2(n-1)}, \ n=1,2,3,4; \notag\\ \lambda_{k,n,j}=s^4_{k,n,j}, \ j=5,6; \ n=1,2,3,4.\notag

\end{gather}

\end{theorem}



Коэффициенты разложений \eqref{mitr:eq62}--\eqref{mitr:eq64} находятся аналогично получению формул \eqref{mitr:eq52}, \eqref{mitr:eq53}, \eqref{mitr:eq57}.







\AdditionalContent{Конкурирующие интересы}{Конкурирующих интересов не имею.} 

\AdditionalContent{Авторская ответственность}{Я несу полную ответственность за предоставление окончательной версии рукописи в печать. Окончательная версия рукописи мною одобрена.} 

\AdditionalContent{Финансирование}{Исследование выполнялось без финансирования.}





\begin{thebibliography}{99}



\RBibitem{mitr:bib1}

\by Митрохин~С.~И.

\paper О <<расщеплении>> кратных в главном собственных значений многоточечных краевых задач

\jour Изв. вузов. Матем.

\yr 1997

\issue 3

\pages 38--43

\mathnet{http://mi.mathnet.ru/ivm1522}

\mathscinet{http://www.ams.org/mathscinet-getitem?mr=1480768}

\zmath{https://zbmath.org/?q=an:0905.34016}

%\transl

%\jour Russian Math. (Iz. VUZ)

%\yr 1997

%\vol 41

%\issue 3

%\pages 37--42





\RBibitem{mitr:bib2}

\by Винокуров~В.~А., Садовничий~В.~А.

\paper Асимптотика любого порядка собственных значений и собственных функций краевой задачи 

Штурма"--~Лиувилля на отрезке с~суммируемым потенциалом

\jour Дифференц. уравнения

\yr 1998

\vol 34

\issue 10

\pages 1423--1426

\mathnet{http://mi.mathnet.ru/de9799}

\mathscinet{http://www.ams.org/mathscinet-getitem?mr=1713014}

%\transl

%\jour Differ. Equ.

%\yr 1998

%\vol 34

%\issue 10

%\pages 1425--1429





\RBibitem{mitr:bib3}

\by Винокуров~В.~А., Садовничий~В.~А.

\paper Асимптотика любого порядка собственных значений и собственных функций краевой задачи Штурма--Лиувилля на~отрезке с~суммируемым потенциалом

\jour Изв. РАН. Сер. матем.

\yr 2000

\vol 64

\issue 4

\pages 47--108

\mathnet{http://mi.mathnet.ru/izv295}

\crossref{10.4213/im295}

\mathscinet{http://www.ams.org/mathscinet-getitem?mr=1794595}

\zmath{https://zbmath.org/?q=an:1001.34021}

%\transl

%\jour Izv. Math.

%\yr 2000

%\vol 64

%\issue 4

%\pages 695--754

%\crossref{10.1070/im2000v064n04ABEH000295}

%\isi{http://gateway.isiknowledge.com/gateway/Gateway.cgi?GWVersion=2&SrcApp=PARTNER_APP&SrcAuth=LinksAMR&DestLinkType=FullRecord&DestApp=ALL_WOS&KeyUT=000165984800002}

%\scopus{http://www.scopus.com/record/display.url?origin=inward&eid=2-s2.0-33746976591}



\RBibitem{mitr:bib4}

\by Митрохин~С.~И.

\paper Асимптотика собственных значений дифференциального оператора четвёртого порядка с суммируемыми коэффициентами

\jour Вестник Московского университета. Сер. Матем., мех.

\yr 2009

\issue 3

\pages 14--17



\RBibitem{mitr:bib5}

\by  Митрохин~С.~И.

\paper О спектральных свойствах одного дифференциального оператора с~суммируемыми коэффициентами с~запаздывающим аргументом

\jour Уфимск. матем. журн.

\yr 2011

\vol 3

\issue 4

\pages 95--115

\mathnet{http://mi.mathnet.ru/ufa121}

\zmath{https://zbmath.org/?q=an:1249.34182}



\RBibitem{mitr:bib6}

\by Ильин~В.~А.

\paper О~сходимости разложений по собственным функциям в точках разрыва коэффициентов дифференциального оператора

\jour Матем. заметки

\yr 1977

\vol 22

\issue 5

\pages 679--698

\mathnet{http://mi.mathnet.ru/mz8092}

\mathscinet{http://www.ams.org/mathscinet-getitem?mr=499820}

\zmath{https://zbmath.org/?q=an:0364.34011}

%\transl

%\jour Math. Notes

%\yr 1977

%\vol 22

%\issue 5

%\pages 870--882

%\crossref{10.1007/BF01098352}



\RBibitem{mitr:bib7}

\by Будаев~В.~Д. 

\paper О~безусловной базисности на замкнутом интервале систем собственных и присоединенных 

функций оператора второго порядка с~разрывными коэффициентами

\jour Дифференц. уравнения

\yr 1987

\vol 23

\issue 6

\pages 941--952

\mathnet{http://mi.mathnet.ru/de6222}

\mathscinet{http://www.ams.org/mathscinet-getitem?mr=899614}

\zmath{https://zbmath.org/?q=an:0636.34017}



\RBibitem{mitr:bib8}

\by Ильин~В.~А.

\paper Необходимые и достаточные условия базисности Рисса корневых векторов разрывных операторов 

второго порядка

\jour Дифференц. уравнения

\yr 1986

\vol 22

\issue 12

\pages 2059--2071

\mathnet{http://mi.mathnet.ru/de6059}

\mathscinet{http://www.ams.org/mathscinet-getitem?mr=880763}







\RBibitem{mitr:bib9}

\by Митрохин~С.~И.

\paper О некоторых спектральных свойствах дифференциальных операторов второго порядка с разрывной весовой функцией 

\jour Докл. РАН

\yr 1997

\vol 356

\issue 1

\pages 13--15

\mathnet{http://mi.mathnet.ru/dan3659}

\mathscinet{http://www.ams.org/mathscinet-getitem?mr=1492027}



\RBibitem{mitr:bib10}

\by	Гуревич~А.~П., Хромов~А.~П.

\paper Операторы дифференцирования первого и второго порядков со знакопеременной весовой функцией

\jour Матем. заметки

\yr 1994

\vol 56

\issue 1

\pages 3--15

\mathnet{http://mi.mathnet.ru/mz2218}

\mathscinet{http://www.ams.org/mathscinet-getitem?mr=1309814}

\zmath{https://zbmath.org/?q=an:0835.34093}

%\transl

%\jour Math. Notes

%\yr 1994

%\vol 56

%\issue 1

%\pages 653--661

%\crossref{10.1007/BF02110552}

%\isi{http://gateway.isiknowledge.com/gateway/Gateway.cgi?GWVersion=2&SrcApp=PARTNER_APP&SrcAuth=LinksAMR&DestLinkType=FullRecord&DestApp=ALL_WOS&KeyUT=A1994RH68200001}



\RBibitem{mitr:bib11}

\by	Лидский~В.~В., Садовничий~В.~А.

\paper Регуляризованные суммы корней одного класса целых функций

\jour Функц. анализ и его прил.

\yr 1967

\vol 1

\issue 2

\pages 52--59

\mathnet{http://mi.mathnet.ru/faa2816}

\mathscinet{http://www.ams.org/mathscinet-getitem?mr=213520}

\zmath{https://zbmath.org/?q=an:0176.37001}

%\transl

%\jour Funct. Anal. Appl.

%\yr 1967

%\vol 1

%\issue 2

%\pages 133--139

%\crossref{10.1007/BF01076085}



\RBibitem{mitr:bib12}

\by Митрохин~С.~И.

\paper О~спектральных свойствах дифференциальных операторов с~разрывными коэффициентами

\jour Дифференц. уравнения

\yr 1992

\vol 28

\issue 3

\pages 530--532

\mathnet{http://mi.mathnet.ru/de7760}

\mathscinet{http://www.ams.org/mathscinet-getitem?mr=1190396}

\zmath{https://zbmath.org/?q=an:0824.47040}



\RBibitem{mitr:bib13}

\by Лидский~В.~Б., Садовничий~В.~А.

\paper Асимптотические формулы для корней одного класса целых функций

\jour Матем. сб.

\yr 1968

\vol 75(117)

\issue 4

\pages 558--566

\mathnet{http://mi.mathnet.ru/msb4001}

\mathscinet{http://www.ams.org/mathscinet-getitem?mr=223573}

\zmath{https://zbmath.org/?q=an:0176.37002}

%\transl

%\by V.~B.~Lidskii, V.~A.~Sadovnichii

%\paper Asymptotic formulas for the zeros of a class of entire functions

%\jour Math. USSR-Sb.

%\yr 1968

%\vol 4

%\issue 4

%\pages 519--527

%\crossref{10.1070/SM1968v004n04ABEH002812}



\RBibitem{mitr:bib14}

\by Савчук~А.~М.

\paper Регуляризованный след первого порядка оператора Штурма"--~Лиувилля с $\delta$"=потенциалом

\jour УМН

\yr 2000

\vol 55

\issue 6(336)

\pages 155--156

\mathnet{http://mi.mathnet.ru/umn352}

\crossref{10.4213/rm352}

\mathscinet{http://www.ams.org/mathscinet-getitem?mr=1840377}

\zmath{https://zbmath.org/?q=an:1015.34075}

\adsnasa{http://adsabs.harvard.edu/cgi-bin/bib_query?2000RuMaS..55.1168S}

%\transl

%\jour Russian Math. Surveys

%\yr 2000

%\vol 55

%\issue 6

%\pages 1168--1169

%\crossref{10.1070/rm2000v055n06ABEH000352}

%\isi{http://gateway.isiknowledge.com/gateway/Gateway.cgi?GWVersion=2&SrcApp=PARTNER_APP&SrcAuth=LinksAMR&DestLinkType=FullRecord&DestApp=ALL_WOS&KeyUT=000169375800019}

%\scopus{http://www.scopus.com/record/display.url?origin=inward&eid=2-s2.0-0034561354}





\RBibitem{mitr:bib15}

\by Савчук~А.~М., Шкаликов~А.~А.

\paper Операторы Штурма"--~Лиувилля с сингулярными потенциалами

\jour Матем. заметки

\yr 1999

\vol 66

\issue 6

\pages 897--912

\mathnet{http://mi.mathnet.ru/mz1234}

\crossref{10.4213/mzm1234}

\mathscinet{http://www.ams.org/mathscinet-getitem?mr=1756602}

\zmath{https://zbmath.org/?q=an:0968.34072}

\elib{http://elibrary.ru/item.asp?id=13650063}

%\transl

%\jour Math. Notes

%\yr 1999

%\vol 66

%\issue 6

%\pages 741--753

%\crossref{10.1007/BF02674332}

%\isi{http://gateway.isiknowledge.com/gateway/Gateway.cgi?GWVersion=2&SrcApp=PARTNER_APP&SrcAuth=LinksAMR&DestLinkType=FullRecord&DestApp=ALL_WOS&KeyUT=000086576400027}



\RBibitem{mitr:bib16}

\by Наймарк~М.~А.

\book Линейные дифференциальные операторы

\publaddr М.

\publ Наука

\yr 1969

\totalpages 528

\zmath{https://zbmath.org/?q=an:0193.04101}



\RBibitem{mitr:bib17}

\by Левитан~Б.~М., Саргсян~И.~С. 

\book Введение в спектральную теорию. Самосопряженные обыкновенные дифференциальные операторы

\publaddr  М.

\publ Наука

\yr 1970

\totalpages 672

\zmath{https://zbmath.org/?q=an:0225.47019}



\Bibitem{mitr:bib18}

\by Bellman~R., Cooke~K.~L.

\book Differential-difference equations

\zmath{https://zbmath.org/?q=an:0105.06402}

\serial Mathematics in Science and Engineering

\vol 6

\publaddr New York--London

\publ Academic Press

\totalpages xvi+462 

\yr 1963



\RBibitem{mitr:bib19}

\by Садовничий~В.~А., Любишкин~В.~А.

\paper О~некоторых новых результатах теории регуляризованных следов дифференциальных операторов

\jour Дифференц. уравнения

\yr 1982

\vol 18

\issue 1

\pages 109--116

\mathnet{http://mi.mathnet.ru/de4469}

\mathscinet{http://www.ams.org/mathscinet-getitem?mr=646537}



\RBibitem{mitr:bib20}

\by Садовничий~В.~А., Любишкин~В.~А., Белабасси~Ю.

\paper О регуляризованных суммах корней целой функции одного класса

\jour Докл. АН СССР

\yr 1980

\vol 254

\issue 6

\pages 1346--1348

\zmath{https://zbmath.org/?q=an:0474.30027}



\RBibitem{mitr:bib21}

\by Митрохин~С.~И.

\paper О спектральных свойствах дифференциального оператора Штурма"--~Лиувилля с запаздывающим аргументом

\jour Вестник Московского университета. Сер. Матем., мех.

\yr 2013

\issue 4

\pages 38--42



\RBibitem{mitr:bib22}

\by Митрохин~С.~И.

\paper О спектральных свойствах дифференциальных операторов нечетного порядка с суммируемым потенциалом 

\jour Дифференц. уравнения

\yr 2011

\vol 47

\issue 12

\pages 1808--1811



\end{thebibliography}







\newpage







\makeentitle







\AdditionalContent{Competing interests}{I have no competing interests.} 

\AdditionalContent{Author's Responsibilities}{I take full responsibility for submitting the final manuscript in print. I approved the final version of the manuscript.} 

\AdditionalContent{Funding}{The research has not had any sponsorship.}



\newpage

\begin{thebibliography}{99}



\Bibitem{mitr:bib1:en}

\by Mitrokhin~S.~I.

\paper On the ``splitting'' of multiplicity in principal eigenvalues of multipoint boundary value problems

\mathscinet{http://www.ams.org/mathscinet-getitem?mr=1480768}

\zmath{https://zbmath.org/?q=an:0905.34016}

\jour Russian Math. (Iz. VUZ)

\yr 1997

\vol 41

\issue 3

\pages 37--42





\Bibitem{mitr:bib2:en}

\by Vinokurov~V.~A., Sadovnichii~V.~A.

\paper Arbitrary-order asymptotics of the eigenvalues and eigenfunctions of the Sturm--Liouville boundary value problem on an interval with integrable potential.

\mathscinet{http://www.ams.org/mathscinet-getitem?mr=1713014}

\jour Differ. Equ.

\yr 1998

\vol 34

\issue 10

\pages 1425--1429





\Bibitem{mitr:bib3}

\paper Asymptotics of any order for the eigenvalues and eigenfunctions of the Sturm--Liouville boundary-value problem on a segment with a summable potential

\by Vinokurov~V.~A., Sadovnichii~V.~A.

\jour Izv. Math.

\yr 2000

\vol 64

\issue 4

\pages 695--754

\crossref{10.1070/im2000v064n04ABEH000295}

\isi{http://gateway.isiknowledge.com/gateway/Gateway.cgi?GWVersion=2&SrcApp=PARTNER_APP&SrcAuth=LinksAMR&DestLinkType=FullRecord&DestApp=ALL_WOS&KeyUT=000165984800002}

\scopus{http://www.scopus.com/record/display.url?origin=inward&eid=2-s2.0-33746976591}



\Bibitem{mitr:bib4:en}

\paper  The asymptotics of the eigenvalues of a fourth order differential operator with summable coefficients

\by Mitrokhin~S.~I.

\yr 2009

\jour Moscow Univ. Math. Bull. 

\vol 64

\issue 3

\pages 102--104

\crossref{10.3103/S0027132209030024}

\scopus{2-s2.0-84859613451}





\Bibitem{mitr:bib5:en}

\lang In Russian

\by Mitrokhin~S.~I.

\paper On spectral properties of a~differential operator with summable coefficients with a~retarded argument

\jour Ufimsk. Mat. Zh.

\yr 2011

\vol 3

\issue 4

\pages 95--115





\Bibitem{mitr:bib6:en}

\paper Convergence of eigenfunction expansions at points of discontinuity of the coefficients of a differential operator

\by Il'in~V.~A.

\jour Math. Notes

\yr 1977

\vol 22

\issue 5

\pages 870--882

\crossref{10.1007/BF01098352}



\Bibitem{mitr:bib7:en}

\lang In Russian

\by Budaev~V.~D.

\paper The property of being an unconditional basis on a closed interval, for systems of eigen- and associated functions of a second-order operator with discontinuous coefficients

\jour Differ. Uravn.

\yr 1987

\vol 23

\issue 6

\pages 941--952







\Bibitem{mitr:bib8:en}

\lang In Russian

\by Il'in~V.~A.

\paper Necessary and sufficient conditions for being a Riesz basis of root vectors of second-order discontinuous operators

\jour Differ. Uravn.

\yr 1986

\vol 22

\issue 12

\pages 2059--2071







\Bibitem{mitr:bib9:en}

\lang In Russian

\by Mitrokhin~S.~I.

\paper On some spectral properties of second-order differential operators with a discontinuous positive weight function

\jour Dokl. Akad. Nauk 

\yr 1997

\vol 356

\issue 1

\pages 13--15





\Bibitem{mitr:bib10:en}

\paper First and second order differentiation operators with weight functions of variable sign

\by Gurevich~A.~P.,  Khromov~A.~P.

\jour Math. Notes

\yr 1994

\vol 56

\issue 1

\pages 653--661

\crossref{10.1007/BF02110552}

\mathscinet{http://www.ams.org/mathscinet-getitem?mr=1309814}

\zmath{https://zbmath.org/?q=an:0835.34093}

\isi{http://gateway.isiknowledge.com/gateway/Gateway.cgi?GWVersion=2&SrcApp=PARTNER_APP&SrcAuth=LinksAMR&DestLinkType=FullRecord&DestApp=ALL_WOS&KeyUT=A1994RH68200001}



\Bibitem{mitr:bib11:en}

\paper Regularized sums of zeros of a class of entire functions

\by  Lidskii~V.~B., Sadovnichii~V.~A.

\mathscinet{http://www.ams.org/mathscinet-getitem?mr=213520}

\zmath{https://zbmath.org/?q=an:0176.37001}

\jour Funct. Anal. Appl.

\yr 1967

\vol 1

\issue 2

\pages 133--139

\crossref{10.1007/BF01076085}



\Bibitem{mitr:bib12:en}

\lang In Russian

\by Mitrokhin~S.~I.

\paper Spectral properties of differential operators with discontinuous coefficients

\jour Differ. Uravn.

\yr 1992

\vol 28

\issue 3

\pages 530--532





\Bibitem{mitr:bib13:en}

\mathscinet{http://www.ams.org/mathscinet-getitem?mr=223573}

\zmath{https://zbmath.org/?q=an:0176.37002}

\by Lidskii~V.~B., Sadovnichii~V.~A.

\paper Asymptotic formulas for the zeros of a class of entire functions

\jour Math. USSR-Sb.

\yr 1968

\vol 4

\issue 4

\pages 519--527

\crossref{10.1070/SM1968v004n04ABEH002812}



\Bibitem{mitr:bib14:en}

\paper First-order regularised trace of the Sturm--Liouville operator with $\delta$-potential

\by Savchuk~A.~M.

\mathscinet{http://www.ams.org/mathscinet-getitem?mr=1840377}

\zmath{https://zbmath.org/?q=an:1015.34075}

\adsnasa{http://adsabs.harvard.edu/cgi-bin/bib_query?2000RuMaS..55.1168S}

\jour Russian Math. Surveys

\yr 2000

\vol 55

\issue 6

\pages 1168--1169

\crossref{10.1070/rm2000v055n06ABEH000352}

\isi{http://gateway.isiknowledge.com/gateway/Gateway.cgi?GWVersion=2&SrcApp=PARTNER_APP&SrcAuth=LinksAMR&DestLinkType=FullRecord&DestApp=ALL_WOS&KeyUT=000169375800019}

\scopus{http://www.scopus.com/record/display.url?origin=inward&eid=2-s2.0-0034561354}





\Bibitem{mitr:bib15:en}

\paper Sturm--Liouville operators with singular potentials

\by Savchuk~A.~M., Shkalikov~A.~A.

\mathscinet{http://www.ams.org/mathscinet-getitem?mr=1756602}

\zmath{https://zbmath.org/?q=an:0968.34072}

\jour Math. Notes

\yr 1999

\vol 66

\issue 6

\pages 741--753

\crossref{10.1007/BF02674332}

\isi{http://gateway.isiknowledge.com/gateway/Gateway.cgi?GWVersion=2&SrcApp=PARTNER_APP&SrcAuth=LinksAMR&DestLinkType=FullRecord&DestApp=ALL_WOS&KeyUT=000086576400027}



\Bibitem{mitr:bib16:en}

\lang In Russian

\by Naimark~M.~A.

\book Lineinye differentsial'nye operatory \rm [Linear Differential Operators]

\publaddr Moscow

\publ Nauka

\yr 1969

\totalpages 528





\Bibitem{mitr:bib17:en}

\zmath{https://zbmath.org/?q=an:0302.47036}

\by Levitan~B.~M., Sargsian~I.~S.

\book Introduction to spectral theory. Selfadjoint ordinary differential operators

\serial Translations of Mathematical Monographs

\vol 39

\publaddr Providence, R.I.

\publ American Mathematical Society 

\totalpages xi+525 

\yr 1975



\Bibitem{mitr:bib18:en}

\by Bellman~R., Cooke~K.~L.

\book Differential-difference equations

\zmath{https://zbmath.org/?q=an:0105.06402}

\serial Mathematics in Science and Engineering

\vol 6

\publaddr New York--London

\publ Academic Press

\totalpages xvi+462 

\yr 1963



\Bibitem{mitr:bib19:en}

\lang In Russian

\by Sadovnichii~V.~A., Lyubishkin~V.~A.

\paper Some new results of the theory of regularized traces of differential operators

\jour Differ. Uravn.

\yr 1982

\vol 18

\issue 1

\pages 109--116



\Bibitem{mitr:bib20:en}

\by Sadovnichii~V.~A., Lyubishkin~V.~A., Belabbasi~Yu.

\paper On regularized sums of roots of an entire function of a certain class

\zmath{https://zbmath.org/?q=an:0474.30027}

\jour Sov. Math., Dokl. 

\vol 22

\pages 613--616 

\yr 1980



\Bibitem{mitr:bib21:en}

\jour Moscow Univ. Math. Bull. 

\crossref{10.3103/S0027132213040062}

\scopus{2-s2.0-84883628938}

\vol 68

\issue 4

\yr 2013

\pages 198--201

\paper  Spectral properties of a Sturm--Liouville type differential operator with a retarding argument

\by Mitrokhin~S.~I. 



\Bibitem{mitr:bib22:en}

\jour Differ. Equ.

\yr 2011

\crossref{10.1134/S0012266111120123}

\scopus{2-s2.0-84856539436}

\vol 47

\issue 12

\yr 2011

\pages 1833--1836

\paper On the spectral properties of odd-order differential operators with integrable potential

\by Mitrokhin~S.~I. 



\end{thebibliography}







\renewcommand{\baselinestretch}{0.9}



%\end{document}

